\documentclass[11pt,a4paper]{article}

\usepackage[utf8]{inputenc}
\usepackage[T1]{fontenc}
\usepackage{lmodern}
\usepackage[spanish,es-noshorthands]{babel}
\usepackage[a4paper,margin=2.4cm]{geometry}
\usepackage{amsmath,amssymb}
\usepackage{booktabs}
\usepackage{graphicx}
\usepackage{subcaption}
\usepackage[font=small,labelfont=bf]{caption}
\usepackage{xcolor}
\usepackage{enumitem}
\usepackage{float}
\usepackage{microtype}
\usepackage{fancyhdr}
\usepackage{needspace}
\usepackage[hidelinks]{hyperref}

\addto\captionsspanish{\renewcommand{\tablename}{Tabla}\renewcommand{\contentsname}{Contenido}}

\setlist{itemsep=2pt,topsep=4pt}
\setlength{\parskip}{4pt}
\setlength{\parindent}{0pt}
\renewcommand{\arraystretch}{1.12}

\pagestyle{fancy}
\fancyhf{}
\fancyhead[L]{\small Taller 3: Ordenamiento y B\'usqueda}
\fancyhead[R]{\small Universidad del Quind\'io}
\fancyfoot[C]{\small\thepage}
\renewcommand{\headrulewidth}{0.3pt}

% Texto que el grupo debe completar o confirmar antes de entregar.
\newcommand{\pendiente}[1]{\textcolor{red}{\textbf{[#1]}}}
\newcommand{\bigO}[1]{\ensuremath{O(#1)}}
\newcommand{\code}[1]{\texttt{#1}}

% Dos gr\'aficas (equipo A y equipo B), una debajo de la otra, a ancho completo.
% #1 nombre base del PNG, #2 ancho relativo, #3 pie de figura, #4 etiqueta
\newcommand{\figpair}[4]{%
  \begin{figure}[htbp]
    \centering
    \begin{subfigure}[t]{\linewidth}
      \centering
      \includegraphics[width=#2\linewidth]{figs/A_#1.png}
      \caption{Equipo A.}
    \end{subfigure}\par\medskip
    \begin{subfigure}[t]{\linewidth}
      \centering
      \includegraphics[width=#2\linewidth]{figs/B_#1.png}
      \caption{Equipo B.}
    \end{subfigure}
    \caption{#3}\label{#4}
  \end{figure}}

\begin{document}

% ====================================================================================== PORTADA
\begin{titlepage}
  \centering
  \vspace*{2.2cm}
  {\large Universidad del Quind\'io\par}
  \vspace{2.2cm}
  {\Huge\bfseries Taller 3\par}
  \vspace{0.4cm}
  {\LARGE\bfseries Ordenamiento y B\'usqueda\par}
  \vspace{0.6cm}
  {\large Comparaci\'on experimental de algoritmos de ordenamiento y b\'usqueda\\en Java y Python\par}
  \vspace{2.6cm}
  {\large Santiago Orozco Zuluaga\par}
  {\large Jhony Alexander Perea Perea\par}
  \vspace{1.6cm}
  {\large Curso: An\'alisis de Algoritmos \pendiente{confirmar nombre del curso y docente}\par}
  \vspace{0.6cm}
  {\large Repositorio:\\ \url{https://github.com/Shuh3n/taller3-ordenamiento-busqueda}\par}
  \vfill
  {\large Octubre de 2026\par}
\end{titlepage}

\tableofcontents
\clearpage

% ================================================================================== PRESENTACI\'ON
\section{Presentaci\'on}

Este informe documenta un estudio experimental de seis algoritmos de ordenamiento (Insertion Sort,
Dual-Pivot QuickSort, Merge Sort, Heap Sort, Radix Sort y Bucket Sort) y seis algoritmos de b\'usqueda
(lineal, binaria, ternaria, por saltos, por interpolaci\'on y exponencial). Cada algoritmo se implement\'o
con la misma l\'ogica en dos lenguajes: \textbf{Java} (compilado, con compilaci\'on JIT) y \textbf{Python}
(interpretado), para un total de 24 implementaciones. El objetivo es medir su desempe\~no, verificar si el
comportamiento observado corresponde con la complejidad te\'orica y explicar, a partir de los datos
obtenidos, las diferencias entre algoritmos, entre tipos de entrada y entre lenguajes.

\paragraph{C\'omo se hizo.} Los datos de prueba se generaron una sola vez con una semilla fija y se guardaron
en archivos de texto plano; todos los programas los leen desde esos archivos, de modo que cada algoritmo se
eval\'ua sobre exactamente los mismos valores. Cada medici\'on se repiti\'o cinco veces y se reporta el promedio
y la desviaci\'on est\'andar. Los resultados se almacenaron en archivos CSV y las gr\'aficas y tablas de este
documento se construyeron directamente a partir de ellos.

\paragraph{Dos equipos de c\'omputo.} Todo el flujo experimental se ejecut\'o completo en \textbf{dos
computadores distintos} (equipo A y equipo B), siguiendo el procedimiento de la secci\'on~\ref{sec:metodologia}.
Cada conjunto de resultados sale de un \'unico equipo, para que sus tiempos sean comparables entre s\'i. Las gr\'aficas de cada punto se presentan para ambos
equipos, y la secci\'on~\ref{sec:equipos} compara los dos conjuntos de resultados. Que dos m\'aquinas
independientes lleguen a las mismas conclusiones cualitativas es, adem\'as, una verificaci\'on de que los
hallazgos no dependen de un equipo en particular.

\paragraph{Estructura del documento.} La secci\'on~\ref{sec:metodologia} describe la metodolog\'ia y las
condiciones del experimento; las secciones~\ref{sec:p1}, \ref{sec:p2} y~\ref{sec:p3} desarrollan los Puntos~1
(ordenamiento), 2 (b\'usqueda) y 3 (an\'alisis comparativo) del taller; la secci\'on~\ref{sec:equipos} compara
los dos equipos; despu\'es se declara el uso de inteligencia artificial y se presentan las conclusiones. Las
tablas completas de resultados est\'an en el Anexo~\ref{anexo}.

% ================================================================================== METODOLOG\'IA
\section{Metodolog\'ia y condiciones del experimento}\label{sec:metodologia}

\paragraph{Implementaci\'on.} Las implementaciones de Java y de Python siguen exactamente la misma l\'ogica y
se adaptaron de los art\'iculos de GeeksforGeeks citados en la Tabla~\ref{tab:fuentes}. Los ordenamientos
trabajan \emph{en sitio} sobre el arreglo recibido y las b\'usquedas devuelven un \'indice o $-1$. Dos
adaptaciones merecen mencionarse:
\begin{itemize}
  \item \textbf{Dual-Pivot QuickSort:} los pivotes son el primer y el \'ultimo elemento del intervalo (si el
  primero es mayor, se intercambian), tal como pide el taller. No se usa ninguna estrategia de pivote
  aleatorio ni de mediana.
  \item \textbf{Bucket Sort:} se usan $k=n$ cubetas de igual ancho sobre el intervalo $[\min,\max]$ del
  arreglo. El valor $v$ se asigna a la cubeta $\left\lfloor \dfrac{(v-\min)\cdot k}{\max-\min+1}\right\rfloor$.
  Las cubetas se almacenan de forma contigua (conteo y suma de prefijos) en lugar de una lista por cubeta, y
  cada cubeta se ordena con Insertion Sort, como en GeeksforGeeks. Radix Sort es LSD en base 10, con un
  Counting Sort estable por d\'igito (8 pasadas para n\'umeros de 8 d\'igitos).
\end{itemize}

\begin{table}[htbp]
\centering
\caption{Fuentes (GeeksforGeeks) utilizadas para cada algoritmo.}\label{tab:fuentes}
{\footnotesize
\begin{tabular}{ll}
\toprule
Algoritmo & Fuente \\
\midrule
Insertion Sort & \url{https://www.geeksforgeeks.org/insertion-sort-algorithm/} \\
Dual-Pivot QuickSort & \url{https://www.geeksforgeeks.org/dual-pivot-quicksort/} \\
Merge Sort & \url{https://www.geeksforgeeks.org/merge-sort/} \\
Heap Sort & \url{https://www.geeksforgeeks.org/heap-sort/} \\
Radix Sort & \url{https://www.geeksforgeeks.org/radix-sort/} \\
Bucket Sort & \url{https://www.geeksforgeeks.org/bucket-sort-2/} \\
B\'usqueda lineal & \url{https://www.geeksforgeeks.org/linear-search/} \\
B\'usqueda binaria & \url{https://www.geeksforgeeks.org/binary-search/} \\
B\'usqueda ternaria & \url{https://www.geeksforgeeks.org/ternary-search/} \\
B\'usqueda por saltos & \url{https://www.geeksforgeeks.org/jump-search/} \\
B\'usqueda por interpolaci\'on & \url{https://www.geeksforgeeks.org/interpolation-search/} \\
B\'usqueda exponencial & \url{https://www.geeksforgeeks.org/exponential-search/} \\
\bottomrule
\end{tabular}}
\end{table}

\paragraph{Datos de prueba.} Se generaron una sola vez con el programa \code{tools/generate\_data.py} (semilla
fija 42) y se guardaron en archivos de texto plano, un entero por l\'inea. Para cada tama\~no
$n\in\{1.000;\,10.000;\,100.000;\,1.000.000\}$ hay un arreglo de enteros aleatorios de 8 d\'igitos (entre
$10.000.000$ y $99.999.999$), sus versiones ordenadas ascendente y descendentemente, y un archivo de 1.000
claves de b\'usqueda (500 presentes en el arreglo y 500 ausentes, mezcladas). Los programas leen esos
archivos antes de medir; la lectura y escritura de archivos nunca entra en el tiempo medido.

\paragraph{Medici\'on del tiempo.} Se mide \'unicamente la ejecuci\'on del algoritmo con
\code{System.nanoTime()} en Java y \code{time.perf\_counter\_ns()} en Python. Cada medici\'on se repite
5~veces, \emph{cada una sobre una copia nueva del arreglo original} (la copia no se cronometra, de modo que
la segunda repetici\'on no ordena un arreglo ya ordenado), y se reportan el promedio y la desviaci\'on
est\'andar muestral. En Java se ejecuta antes un calentamiento (10~ejecuciones sobre 10.000 elementos por
algoritmo de ordenamiento y 300~pasadas del conjunto de claves por algoritmo de b\'usqueda) que no entra en el
promedio. Despu\'es de cada ejecuci\'on, fuera del tiempo medido, se verifica que el arreglo qued\'o ordenado
(y con la misma suma de elementos) o que cada b\'usqueda devolvi\'o el resultado correcto.

\paragraph{Mediciones de b\'usqueda.} Una sola b\'usqueda dura tan poco que medirla de forma aislada estar\'ia
dominado por el ruido del reloj. Por eso se busca el conjunto de 1.000 claves en una sola pasada cronometrada y
se reporta el \emph{tiempo promedio por b\'usqueda}. Para el efecto de la posici\'on de la clave, se busca
1.000~veces una misma clave ubicada en el \'indice $0$ (inicio), $n/2$ (centro) o $n-1$ (final) del arreglo de
$n=1.000.000$.

\paragraph{Valores estimados y fallos.} Antes de ejecutar un tama\~no $n$, el programa estima su tiempo
escalando el \'ultimo promedio medido con el modelo de costo te\'orico del algoritmo:
\begin{equation}
  t(n)\;=\;t(n_{\text{prev}})\cdot\frac{f(n)}{f(n_{\text{prev}})},
  \qquad f(n)\in\{n,\; n\log_2 n,\; n^2\}.\label{eq:estimacion}
\end{equation}
Si la estimaci\'on supera 300~s por ejecuci\'on, el algoritmo no se ejecuta y el valor se registra como
\emph{estimado}. Si un algoritmo falla (desbordamiento de pila en Java, \code{RecursionError} en Python), el
error se registra sin ocultarlo y los tama\~nos mayores se estiman con el mismo escalamiento. En las gr\'aficas
de l\'ineas los valores estimados se dibujan con l\'inea punteada y marcador hueco; en las de barras, con barra
rayada y el prefijo <<$\sim$>>; en las tablas, con <<$\sim$>> y <<(est.)>>. Todas las gr\'aficas usan escala
logar\'itmica en el eje del tiempo. El registro completo de estimados y fallos de cada equipo est\'a en las
Tablas~\ref{tab:issuesA} y~\ref{tab:issuesB} del anexo.

\paragraph{Reproducci\'on.} El c\'odigo fuente, los resultados y el programa que genera los datos con semilla fija
est\'an en \url{https://github.com/Shuh3n/taller3-ordenamiento-busqueda}. Las tablas de este documento se
generan con \code{tools/make\_latex.py} a partir de \code{results/} (equipo A) y \code{resultsv2/} (equipo B).

\paragraph{Formato de resultados.} Cada fila de los CSV contiene: lenguaje, algoritmo, variante de entrada (o
escenario de b\'usqueda), $n$, el tiempo de cada una de las cinco repeticiones, promedio, desviaci\'on
est\'andar, si el valor fue medido o estimado, y el estado de la ejecuci\'on.

\paragraph{Entorno de pruebas.}\label{par:entorno} La Tabla~\ref{tab:entorno} resume los dos equipos. Las
corridas se hicieron con los benchmarks ejecutados uno tras otro (nunca en paralelo)
El equipo~B estuvo conectado a la corriente \pendiente{confirmar lo mismo para el equipo A y que no hab\'ia otras aplicaciones pesadas}.

\begin{table}[htbp]
\centering
\caption{Entorno de pruebas de cada equipo.}\label{tab:entorno}
{\small
\begin{tabular}{lll}
\toprule
 & Equipo A (\code{results/}) & Equipo B (\code{resultsv2/}) \\
\midrule
Equipo & MacBook Air (M4) & ASUS Vivobook X1605ZA \\
Sistema operativo & macOS 26.6.2 & Fedora Linux 44 (Workstation) \\
Procesador & Apple M4 (arm64) & Intel Core i5-12500H, 12.\textsuperscript{a} gen., 16 hilos \\
Memoria RAM & 16 GB & 16 GB \\
Almacenamiento & SSD de 256 GB & SSD M.2 de 1 TB \\
Java (JIT) & 27.0.1.0.2 & \pendiente{completar versi\'on} \\
Python & 3.15.1 & \pendiente{completar versi\'on} \\
\bottomrule
\end{tabular}}
\end{table}

% ============================================================================== PUNTO 1
\section{Punto 1: algoritmos de ordenamiento}\label{sec:p1}

La Tabla~\ref{tab:complejidad} resume la complejidad te\'orica investigada para cada algoritmo, su estabilidad y
si trabaja \emph{in-place}.

\begin{table}[htbp]
\centering
\caption{Complejidad te\'orica, estabilidad y uso de memoria ($n$ elementos, $d=8$ d\'igitos, $b=10$ base,
$k$ cubetas). \emph{Estable}: conserva el orden relativo de claves iguales. \emph{In-place}: usa memoria
auxiliar $O(1)$ (o $O(\log n)$ de pila).}\label{tab:complejidad}
\resizebox{\linewidth}{!}{\footnotesize
\begin{tabular}{lccccc}
\toprule
Algoritmo & Mejor caso & Caso promedio & Peor caso & Estable & In-place \\
\midrule
Insertion Sort & $O(n)$ & $O(n^2)$ & $O(n^2)$ & S\'i & S\'i \\
Dual-Pivot QuickSort & $O(n\log n)$ & $O(n\log n)$ & $O(n^2)$ & No & S\'i (pila $O(\log n)$ prom.) \\
Merge Sort & $O(n\log n)$ & $O(n\log n)$ & $O(n\log n)$ & S\'i & No ($O(n)$ auxiliar) \\
Heap Sort & $O(n\log n)$ & $O(n\log n)$ & $O(n\log n)$ & No & S\'i \\
Radix Sort (LSD, base 10) & $O(d(n+b))$ & $O(d(n+b))$ & $O(d(n+b))$ & S\'i & No ($O(n+b)$ auxiliar) \\
Bucket Sort & $O(n+k)$ & $O(n+k)$ & $O(n^2)$ & S\'i (con inserci\'on estable) & No ($O(n+k)$ auxiliar) \\
\bottomrule
\end{tabular}}
\end{table}

Dos observaciones sobre esta tabla, relevantes para lo que sigue. Primero, con pivotes fijos en los
extremos, el peor caso de Dual-Pivot QuickSort aparece justamente con una entrada ya ordenada: el pivote
izquierdo es el m\'inimo y el derecho el m\'aximo del intervalo, por lo que cada partici\'on solo descarta dos
elementos y la recursi\'on llega a una profundidad cercana a $n/2$. Segundo, el peor caso de Bucket Sort ocurre
cuando casi todos los elementos caen en una misma cubeta, pues entonces el Insertion Sort interno trabaja sobre
casi todo el arreglo.

% ---------------------------------------------------------------------------- 1.1
\subsection{Crecimiento del tiempo con $n$}

Las Figuras~\ref{fig:p1growthJava} y~\ref{fig:p1growthPython} muestran el tiempo contra $n$ con entrada
aleatoria, una l\'inea por algoritmo y escala logar\'itmica en ambos ejes. En una gr\'afica log-log, un
algoritmo $O(n^2)$ se ve como una recta de pendiente $\approx2$ y uno $O(n)$ u $O(n\log n)$ como una recta de
pendiente $\approx1$.

\figpair{p1_growth_java}{0.82}{Ordenamiento en Java, entrada aleatoria (log-log). Medido: l\'inea continua;
estimado: l\'inea punteada y marcador hueco.}{fig:p1growthJava}
\figpair{p1_growth_python}{0.82}{Ordenamiento en Python, entrada aleatoria (log-log). En Python, Insertion Sort
con $n=1.000.000$ es un valor estimado.}{fig:p1growthPython}

La Tabla~\ref{tab:growth} muestra, para cada algoritmo, \emph{cu\'antas veces aument\'o el tiempo medio al
multiplicar $n$ por 10}, junto al factor que predice la complejidad te\'orica: $\times100$ para $O(n^2)$,
$\times10$ para $O(n)$ y $\times13{,}3$, $\times12{,}5$ y $\times12$ para $O(n\log n)$ en cada uno de los tres
saltos.

\begin{table}[htbp]
\centering
\caption{Factor de crecimiento del tiempo medio al multiplicar $n$ por 10 (entrada aleatoria) frente al factor
que predice la complejidad te\'orica. <<$\sim$>> indica que alguno de los extremos es estimado.}\label{tab:growth}
\textbf{Equipo A}\par\smallskip
\textbf{Java}\par\smallskip
{\footnotesize
\begin{tabular}{lccc}
\toprule
Algoritmo & 1.000 $\to$ 10.000 & 10.000 $\to$ 100.000 & 100.000 $\to$ 1.000.000 \\
\midrule
Insertion Sort & $\times$74,2 (te\'or. $\times$100) & $\times$103 (te\'or. $\times$100) & $\times$102 (te\'or. $\times$100) \\
Dual-Pivot QuickSort & $\times$9,36 (te\'or. $\times$13,3) & $\times$11,2 (te\'or. $\times$12,5) & $\times$11,1 (te\'or. $\times$12) \\
Merge Sort & $\times$11,5 (te\'or. $\times$13,3) & $\times$12 (te\'or. $\times$12,5) & $\times$12 (te\'or. $\times$12) \\
Heap Sort & $\times$2,72 (te\'or. $\times$13,3) & $\times$11,9 (te\'or. $\times$12,5) & $\times$12,6 (te\'or. $\times$12) \\
Radix Sort & $\times$7,63 (te\'or. $\times$10) & $\times$9,64 (te\'or. $\times$10) & $\times$9,42 (te\'or. $\times$10) \\
Bucket Sort & $\times$6,36 (te\'or. $\times$10) & $\times$11,7 (te\'or. $\times$10) & $\times$4,63 (te\'or. $\times$10) \\
\bottomrule
\end{tabular}}\par
\medskip
\textbf{Python}\par\smallskip
{\footnotesize
\begin{tabular}{lccc}
\toprule
Algoritmo & 1.000 $\to$ 10.000 & 10.000 $\to$ 100.000 & 100.000 $\to$ 1.000.000 \\
\midrule
Insertion Sort & $\times$123 (te\'or. $\times$100) & $\times$113 (te\'or. $\times$100) & $\sim$$\times$100 (te\'or. $\times$100) \\
Dual-Pivot QuickSort & $\times$14,9 (te\'or. $\times$13,3) & $\times$12,2 (te\'or. $\times$12,5) & $\times$13,1 (te\'or. $\times$12) \\
Merge Sort & $\times$13,5 (te\'or. $\times$13,3) & $\times$12,2 (te\'or. $\times$12,5) & $\times$12,7 (te\'or. $\times$12) \\
Heap Sort & $\times$14,9 (te\'or. $\times$13,3) & $\times$12,8 (te\'or. $\times$12,5) & $\times$15,6 (te\'or. $\times$12) \\
Radix Sort & $\times$11 (te\'or. $\times$10) & $\times$9,63 (te\'or. $\times$10) & $\times$19,3 (te\'or. $\times$10) \\
Bucket Sort & $\times$12,8 (te\'or. $\times$10) & $\times$11,1 (te\'or. $\times$10) & $\times$18,5 (te\'or. $\times$10) \\
\bottomrule
\end{tabular}}\par
\medskip

\textbf{Equipo B}\par\smallskip
\textbf{Java}\par\smallskip
{\footnotesize
\begin{tabular}{lccc}
\toprule
Algoritmo & 1.000 $\to$ 10.000 & 10.000 $\to$ 100.000 & 100.000 $\to$ 1.000.000 \\
\midrule
Insertion Sort & $\times$70,2 (te\'or. $\times$100) & $\times$107 (te\'or. $\times$100) & $\times$149 (te\'or. $\times$100) \\
Dual-Pivot QuickSort & $\times$11,7 (te\'or. $\times$13,3) & $\times$11,1 (te\'or. $\times$12,5) & $\times$11,6 (te\'or. $\times$12) \\
Merge Sort & $\times$10,8 (te\'or. $\times$13,3) & $\times$12,3 (te\'or. $\times$12,5) & $\times$12,4 (te\'or. $\times$12) \\
Heap Sort & $\times$11 (te\'or. $\times$13,3) & $\times$11,5 (te\'or. $\times$12,5) & $\times$12,1 (te\'or. $\times$12) \\
Radix Sort & $\times$6,49 (te\'or. $\times$10) & $\times$9,55 (te\'or. $\times$10) & $\times$10,9 (te\'or. $\times$10) \\
Bucket Sort & $\times$4,89 (te\'or. $\times$10) & $\times$8,34 (te\'or. $\times$10) & $\times$17,4 (te\'or. $\times$10) \\
\bottomrule
\end{tabular}}\par
\medskip
\textbf{Python}\par\smallskip
{\footnotesize
\begin{tabular}{lccc}
\toprule
Algoritmo & 1.000 $\to$ 10.000 & 10.000 $\to$ 100.000 & 100.000 $\to$ 1.000.000 \\
\midrule
Insertion Sort & $\times$104 (te\'or. $\times$100) & $\times$136 (te\'or. $\times$100) & $\sim$$\times$100 (te\'or. $\times$100) \\
Dual-Pivot QuickSort & $\times$14,3 (te\'or. $\times$13,3) & $\times$12,5 (te\'or. $\times$12,5) & $\times$13,6 (te\'or. $\times$12) \\
Merge Sort & $\times$13 (te\'or. $\times$13,3) & $\times$12,5 (te\'or. $\times$12,5) & $\times$13,5 (te\'or. $\times$12) \\
Heap Sort & $\times$14,5 (te\'or. $\times$13,3) & $\times$13,2 (te\'or. $\times$12,5) & $\times$13,7 (te\'or. $\times$12) \\
Radix Sort & $\times$12,5 (te\'or. $\times$10) & $\times$11 (te\'or. $\times$10) & $\times$17,4 (te\'or. $\times$10) \\
Bucket Sort & $\times$10,2 (te\'or. $\times$10) & $\times$12,4 (te\'or. $\times$10) & $\times$14,6 (te\'or. $\times$10) \\
\bottomrule
\end{tabular}}\par
\medskip

\end{table}

\Needspace{9\baselineskip}\paragraph{An\'alisis.}
\begin{itemize}
  \item \textbf{Insertion Sort} es el \'unico algoritmo cuadr\'atico y se nota: en Java, el tiempo crece
  $\times102$ y $\times103$ en los dos \'ultimos saltos del equipo~A (te\'orico $\times100$), y $\times107$ y
  $\times149$ en el equipo~B. Con $n=1.000.000$ tarda 36,7~s en A y 55,3~s en B, frente a 51,7~ms y 76,9~ms de
  Dual-Pivot QuickSort: ordena el mismo arreglo 700 veces m\'as lento. El $\times149$ del equipo~B (mayor que
  100) y los $\times113$ y $\times136$ que se miden en Python entre $10^4$ y $10^5$ sugieren que, al crecer el
  arreglo, aparece un costo adicional por jerarqu\'ia de memoria (el arreglo de $10^6$ enteros deja de caber en
  las cach\'es m\'as r\'apidas); es una explicaci\'on plausible, no la verificamos con un perfilador.
  \item \textbf{Merge Sort, Heap Sort y Dual-Pivot QuickSort} (todos $O(n\log n)$ sobre datos aleatorios)
  crecen $\times11$ a $\times13$ por cada salto de 10 en Java a partir de $n=10.000$, muy cerca del factor
  te\'orico de $\times12$ a $\times12{,}5$. En Python los factores son algo mayores ($\times12$ a
  $\times15{,}6$). Que el factor medido
  est\'e ligeramente por encima de $\times12$ es esperable: la complejidad te\'orica cuenta operaciones, no
  accesos a memoria, y con arreglos grandes los accesos pierden localidad.
  \item \textbf{Radix Sort y Bucket Sort} deber\'ian crecer $\times10$ ($O(d(n+b))$ con $d=8$ fijo y $O(n+k)$
  con $k=n$). En Java el factor es $\times9{,}4$ a $\times9{,}6$ para Radix Sort en el equipo~A, y $\times10{,}9$
  en B. Bucket Sort es el m\'as irregular ($\times4{,}6$ a $\times17{,}4$): sus accesos a memoria son
  dispersos (cada elemento se escribe en la posici\'on de su cubeta), por lo que su costo por elemento depende
  mucho de si el arreglo cabe en cach\'e. En Python ambos superan $\times10$ en el \'ultimo salto
  ($\times14{,}6$ a $\times19{,}3$), con el mismo tipo de explicaci\'on: una lista de $10^6$ enteros de Python
  son $10^6$ objetos dispersos en memoria, y las escrituras dispersas fallan m\'as en cach\'e.
  \item \textbf{Efecto de $n$ peque\~no.} Los saltos $1.000\to10.000$ de Java se desv\'ian m\'as: Heap Sort
  muestra $\times2{,}7$ en el equipo~A cuando la teor\'ia pide $\times13{,}3$. La causa m\'as probable es que
  la medici\'on de $n=1.000$ (unas decenas de microsegundos) a\'un incluye costos del JIT que el calentamiento
  de 10~ejecuciones no alcanz\'o a absorber: ese primer valor es demasiado alto y el factor sale artificialmente
  bajo.
\end{itemize}

\paragraph{C\'omo se estimaron los casos que no terminaron.} Se us\'o la Ecuaci\'on~\eqref{eq:estimacion}.
Hubo dos grupos (listados en las Tablas~\ref{tab:issuesA} y~\ref{tab:issuesB}):
\begin{enumerate}
  \item \emph{Python, Insertion Sort con $n=1.000.000$} (entrada aleatoria y descendente): se escal\'o con
  $f(n)=n^2$ desde $n=100.000$ ($\times100$). Resultan $\sim$6.890~s ($\approx1{,}9$~h) y $\sim$12.400~s
  ($\approx3{,}4$~h) en el equipo~A, y $\sim$13.700~s ($\approx3{,}8$~h) y $\sim$22.100~s ($\approx6{,}1$~h)
  en el B; todos superan el l\'imite de 5~minutos. Como en las medidas reales de Python el salto
  $10^4\to10^5$ fue $\times113$ (A) y $\times136$ (B), y no exactamente $\times100$, estas estimaciones deben
  leerse como una \emph{cota inferior}.
  \item \emph{Dual-Pivot QuickSort con entrada ascendente o descendente}: el algoritmo fall\'o (Java:
  \code{Stack\-Overflow\-Error} con $n=100.000$; Python: \code{Recursion\-Error} con $n=10.000$, l\'imite
  de recursi\'on 1.000), por lo que los tama\~nos mayores se estimaron con $f(n)=n^2$ desde el \'ultimo tama\~no
  que s\'i termin\'o. La validez de ese modelo se puede comprobar en Java, donde $n=1.000\to10.000$ s\'i se
  midi\'o: el tiempo creci\'o $\times106$ en A y $\times92$ en B, es decir, cuadr\'aticamente.
\end{enumerate}

% ---------------------------------------------------------------------------- 1.2
\subsection{Efecto del orden de entrada (lenguaje compilado: Java)}

Las Figuras~\ref{fig:p1order1k}, \ref{fig:p1order10k} y~\ref{fig:p1order100k} comparan, para
$n=1.000$, $10.000$ y $100.000$, los tres tipos de entrada (aleatoria, ascendente y descendente) en Java. Las
barras rayadas son valores estimados; sobre cada barra se muestra su tiempo.

\figpair{p1_order_1000}{0.92}{Efecto del orden de entrada en Java, $n=1.000$.}{fig:p1order1k}
\figpair{p1_order_10000}{0.92}{Efecto del orden de entrada en Java, $n=10.000$.}{fig:p1order10k}
\figpair{p1_order_100000}{0.92}{Efecto del orden de entrada en Java, $n=100.000$. Las barras rayadas de
Dual-Pivot QuickSort son estimadas: el algoritmo fall\'o con \code{StackOverflowError}.}{fig:p1order100k}

\Needspace{9\baselineskip}\paragraph{An\'alisis.}
\begin{itemize}
  \item \textbf{Insertion Sort cambi\'o radicalmente.} Con entrada ascendente es su mejor caso, $O(n)$: cada
  elemento ya est\'a en su lugar y el ciclo interno no desplaza nada. Con $n=100.000$ tard\'o 53,1~$\mu$s frente
  a 361~ms con entrada aleatoria en el equipo~A (casi 6.800 veces menos), y 76,6~$\mu$s frente a 372~ms en el
  B. Con entrada descendente es su peor caso, $O(n^2)$ completo: cada elemento debe recorrer todos los anteriores
  y tarda $\approx1{,}8$ a $2{,}1$ veces m\'as que con datos aleatorios (en promedio, con datos aleatorios cada
  elemento solo se desplaza la mitad del camino).
  \item \textbf{Dual-Pivot QuickSort fue el m\'as afectado, y en sentido contrario.} Con la entrada ya ordenada
  (en cualquier direcci\'on) los pivotes extremos son el m\'inimo y el m\'aximo, la partici\'on queda
  totalmente desbalanceada y el algoritmo degenera a $O(n^2)$: con $n=1.000$ ya es unas 3~veces m\'as lento que
  con datos aleatorios, con $n=10.000$ es $\times35$ (A) y $\times27$ (B), y con $n=100.000$ la profundidad
  de recursi\'on ($\approx n/2=50.000$ llamadas) agot\'o la pila de Java con un \code{StackOverflowError}. Los
  tiempos de $n=100.000$ de las barras rayadas son estimados con $f(n)=n^2$ ($\sim$1,46~s en A y $\sim$1,64~s
  en B frente a 4,67~ms y 6,61~ms con entrada aleatoria, unas 250 a 310 veces m\'as). Es la raz\'on por la que
  las implementaciones reales eligen el pivote al azar, por mediana de tres o por muestreo.
  \item \textbf{Merge Sort, Heap Sort y Bucket Sort mejoraron moderadamente con entrada ordenada}
  (aproximadamente entre 0,3 y 0,65 veces el tiempo de la entrada aleatoria; la excepci\'on es Heap Sort con
  $n=1.000$ en A, $\approx0{,}16$, afectado por el calentamiento). Su complejidad no cambia ---Merge
  Sort y Heap Sort siempre hacen $\Theta(n\log n)$ operaciones---, pero con datos ordenados las comparaciones
  son muy predecibles (en la mezcla siempre sale primero un mismo lado, el procesador acierta las
  ramificaciones) y los accesos a memoria son m\'as secuenciales. En Bucket Sort, un arreglo ordenado asigna
  los elementos a cubetas consecutivas, lo que vuelve secuencial la escritura que con datos aleatorios es
  dispersa. Son efectos de constantes, no de complejidad.
  \item \textbf{Radix Sort fue el \'unico insensible al orden} (cociente entre 0,65 y 1,29 en todos los
  tama\~nos): siempre hace las mismas 8 pasadas de Counting Sort sobre los d\'igitos, y el costo de cada pasada
  no depende de c\'omo est\'en ordenados los valores, porque nunca compara dos elementos entre s\'i.
  \item \textbf{Peor comportamiento.} Entre los algoritmos que terminaron, el peor tiempo absoluto fue
  Insertion Sort con entrada descendente ($n=1.000.000$: 73,8~s en A y 93,2~s en B). Entre los que \emph{no}
  terminaron, el peor comportamiento fue Dual-Pivot QuickSort con entrada ya ordenada: lo provoc\'o la
  combinaci\'on de pivotes en los extremos (partici\'on de tama\~no $n-2$, recursi\'on de profundidad
  $\approx n/2$). El mismo patr\'on aparece en Python, pero antes: el l\'imite de recursi\'on de 1.000 hace
  fallar el algoritmo con $n=10.000$.
\end{itemize}

% ---------------------------------------------------------------------------- 1.3
\subsection{Compilado frente a interpretado}

La Figura~\ref{fig:p1compiled} compara Java y Python con $n=1.000.000$ y entrada aleatoria, con dos barras por
algoritmo (una por lenguaje).

\figpair{p1_compiled}{0.92}{Java frente a Python, $n=1.000.000$, entrada aleatoria. La barra de Python de
Insertion Sort es un valor estimado.}{fig:p1compiled}

\Needspace{9\baselineskip}\paragraph{An\'alisis.} Java fue m\'as r\'apido en todos los algoritmos, pero por factores muy distintos
(cociente Python/Java en el equipo~A, y entre par\'entesis en el B): Dual-Pivot QuickSort $\times12$ ($\times14$),
Merge Sort $\times13$ ($\times14$), Heap Sort $\times23$ ($\times24$), Bucket Sort $\times54$ ($\times26$),
Radix Sort $\times74$ ($\times38$) e Insertion Sort $\sim\times188$ ($\sim\times248$, con la barra de Python
estimada).
\begin{itemize}
  \item \textbf{Compilaci\'on JIT frente a interpretaci\'on.} El JIT de Java traduce los m\'etodos calientes a
  c\'odigo de m\'aquina y los optimiza; Python ejecuta cada operaci\'on como instrucciones de bytecode que el
  int\'erprete despacha una a una, con comprobaci\'on de tipos en cada operaci\'on.
  \item \textbf{Representaci\'on de los datos.} En Java, \code{int[]} es un bloque contiguo de enteros de 4
  bytes. En Python, una lista guarda \emph{referencias} a objetos \code{int} dispersos por el mont\'on: leer un
  elemento es seguir un puntero, lo que cuesta m\'as y aprovecha peor la cach\'e.
  \item \textbf{Por qu\'e la brecha es mayor en unos algoritmos.} Es mayor donde el trabajo es una cantidad
  enorme de iteraciones simples sin ramificaciones impredecibles: Insertion Sort (ciclo interno con un
  desplazamiento por iteraci\'on), Radix Sort (8 pasadas $\times$ 3 ciclos sobre $10^6$ elementos) y Bucket Sort.
  En Java esos ciclos corren casi a la velocidad de la memoria; en Python cada iteraci\'on cuesta decenas de
  nanosegundos de despacho de bytecode. Es menor en Dual-Pivot QuickSort y Merge Sort: ah\'i parte del tiempo en
  Java se va en costos que no desaparecen en Python (fallos de predicci\'on de ramificaciones en las
  comparaciones, recursi\'on) y en Python algunas operaciones se resuelven en C (en la mezcla, las copias con
  \emph{slicing}), de manera que la distancia relativa se reduce. Esta es nuestra interpretaci\'on de los
  datos; no la medimos de forma aislada.
\end{itemize}

% ============================================================================== PUNTO 2
\section{Punto 2: algoritmos de b\'usqueda}\label{sec:p2}

\begin{table}[htbp]
\centering
\caption{Complejidad te\'orica de los algoritmos de b\'usqueda sobre un arreglo ordenado de $n$ elementos.}\label{tab:busqueda}
\resizebox{\linewidth}{!}{\footnotesize
\begin{tabular}{llll}
\toprule
Algoritmo & Mejor caso & Peor caso / promedio & Supuesto \\
\midrule
Lineal & $O(1)$ & $O(n)$ & Ninguno (no requiere orden) \\
Binaria & $O(1)$ & $O(\log_2 n)$ & Arreglo ordenado \\
Ternaria & $O(1)$ & $O(\log_3 n)$ iteraciones (2 comparaciones c/u) & Arreglo ordenado \\
Por saltos & $O(1)$ & $O(\sqrt n)$ & Arreglo ordenado \\
Interpolaci\'on & $O(1)$ & promedio $O(\log\log n)$; peor caso $O(n)$ & Ordenado y distribuci\'on uniforme \\
Exponencial & $O(1)$ & $O(\log i)$, $i$ = posici\'on de la clave & Arreglo ordenado \\
\bottomrule
\end{tabular}}
\end{table}

Los arreglos de entrada son los ordenados ascendentemente del Punto~1. Cada medici\'on busca el conjunto de
1.000 claves (500 presentes y 500 ausentes), cronometra el total y reporta el promedio por b\'usqueda.

% ---------------------------------------------------------------------------- 2.1
\subsection{Crecimiento del tiempo con $n$}

\figpair{p2_growth_java}{0.82}{Tiempo promedio por b\'usqueda frente a $n$ en Java (log-log).}{fig:p2growthJava}
\figpair{p2_growth_python}{0.82}{Tiempo promedio por b\'usqueda frente a $n$ en Python (log-log).}{fig:p2growthPython}

\begin{table}[htbp]
\centering
\caption{Factor de crecimiento del tiempo promedio por b\'usqueda al multiplicar $n$ por 10, frente al factor
te\'orico ($\times10$ lineal; $\times1{,}33$, $\times1{,}25$, $\times1{,}2$ logar\'itmico; $\times3{,}16$ para
$O(\sqrt n)$; $\times1{,}13$, $\times1{,}09$, $\times1{,}06$ para $O(\log\log n)$).}\label{tab:searchgrowth}
\textbf{Equipo A}\par\smallskip
\textbf{Java}\par\smallskip
{\footnotesize
\begin{tabular}{lccc}
\toprule
Algoritmo & 1.000 $\to$ 10.000 & 10.000 $\to$ 100.000 & 100.000 $\to$ 1.000.000 \\
\midrule
Lineal & $\times$9,59 (te\'or. $\times$10) & $\times$10 (te\'or. $\times$10) & $\times$12,1 (te\'or. $\times$10) \\
Binaria & $\times$2,17 (te\'or. $\times$1,33) & $\times$0,828 (te\'or. $\times$1,25) & $\times$1,6 (te\'or. $\times$1,2) \\
Ternaria & $\times$1,45 (te\'or. $\times$1,33) & $\times$1,36 (te\'or. $\times$1,25) & $\times$1,71 (te\'or. $\times$1,2) \\
Por saltos & $\times$0,748 (te\'or. $\times$3,16) & $\times$2,57 (te\'or. $\times$3,16) & $\times$3,04 (te\'or. $\times$3,16) \\
Interpolaci\'on & $\times$1,42 (te\'or. $\times$1,13) & $\times$1,02 (te\'or. $\times$1,09) & $\times$1,68 (te\'or. $\times$1,06) \\
Exponencial & $\times$1,12 (te\'or. $\times$1,33) & $\times$1,98 (te\'or. $\times$1,25) & $\times$1,54 (te\'or. $\times$1,2) \\
\bottomrule
\end{tabular}}\par
\medskip
\textbf{Python}\par\smallskip
{\footnotesize
\begin{tabular}{lccc}
\toprule
Algoritmo & 1.000 $\to$ 10.000 & 10.000 $\to$ 100.000 & 100.000 $\to$ 1.000.000 \\
\midrule
Lineal & $\times$13,4 (te\'or. $\times$10) & $\times$10,1 (te\'or. $\times$10) & $\times$9,93 (te\'or. $\times$10) \\
Binaria & $\times$1,48 (te\'or. $\times$1,33) & $\times$1,37 (te\'or. $\times$1,25) & $\times$1,32 (te\'or. $\times$1,2) \\
Ternaria & $\times$1,39 (te\'or. $\times$1,33) & $\times$1,34 (te\'or. $\times$1,25) & $\times$1,28 (te\'or. $\times$1,2) \\
Por saltos & $\times$3,48 (te\'or. $\times$3,16) & $\times$3,11 (te\'or. $\times$3,16) & $\times$3,37 (te\'or. $\times$3,16) \\
Interpolaci\'on & $\times$1,29 (te\'or. $\times$1,13) & $\times$1,19 (te\'or. $\times$1,09) & $\times$1,11 (te\'or. $\times$1,06) \\
Exponencial & $\times$1,44 (te\'or. $\times$1,33) & $\times$1,4 (te\'or. $\times$1,25) & $\times$1,29 (te\'or. $\times$1,2) \\
\bottomrule
\end{tabular}}\par
\medskip

\textbf{Equipo B}\par\smallskip
\textbf{Java}\par\smallskip
{\footnotesize
\begin{tabular}{lccc}
\toprule
Algoritmo & 1.000 $\to$ 10.000 & 10.000 $\to$ 100.000 & 100.000 $\to$ 1.000.000 \\
\midrule
Lineal & $\times$11,2 (te\'or. $\times$10) & $\times$9,13 (te\'or. $\times$10) & $\times$9,39 (te\'or. $\times$10) \\
Binaria & $\times$1,34 (te\'or. $\times$1,33) & $\times$1,48 (te\'or. $\times$1,25) & $\times$1,71 (te\'or. $\times$1,2) \\
Ternaria & $\times$1,37 (te\'or. $\times$1,33) & $\times$1,22 (te\'or. $\times$1,25) & $\times$1,82 (te\'or. $\times$1,2) \\
Por saltos & $\times$0,511 (te\'or. $\times$3,16) & $\times$2,39 (te\'or. $\times$3,16) & $\times$3,19 (te\'or. $\times$3,16) \\
Interpolaci\'on & $\times$1,18 (te\'or. $\times$1,13) & $\times$1,44 (te\'or. $\times$1,09) & $\times$1,27 (te\'or. $\times$1,06) \\
Exponencial & $\times$1,33 (te\'or. $\times$1,33) & $\times$1,56 (te\'or. $\times$1,25) & $\times$1,7 (te\'or. $\times$1,2) \\
\bottomrule
\end{tabular}}\par
\medskip
\textbf{Python}\par\smallskip
{\footnotesize
\begin{tabular}{lccc}
\toprule
Algoritmo & 1.000 $\to$ 10.000 & 10.000 $\to$ 100.000 & 100.000 $\to$ 1.000.000 \\
\midrule
Lineal & $\times$10,7 (te\'or. $\times$10) & $\times$9,56 (te\'or. $\times$10) & $\times$9,74 (te\'or. $\times$10) \\
Binaria & $\times$1,44 (te\'or. $\times$1,33) & $\times$1,36 (te\'or. $\times$1,25) & $\times$1,42 (te\'or. $\times$1,2) \\
Ternaria & $\times$1,42 (te\'or. $\times$1,33) & $\times$1,36 (te\'or. $\times$1,25) & $\times$1,38 (te\'or. $\times$1,2) \\
Por saltos & $\times$3,22 (te\'or. $\times$3,16) & $\times$3,12 (te\'or. $\times$3,16) & $\times$2,84 (te\'or. $\times$3,16) \\
Interpolaci\'on & $\times$1,22 (te\'or. $\times$1,13) & $\times$1,19 (te\'or. $\times$1,09) & $\times$1,21 (te\'or. $\times$1,06) \\
Exponencial & $\times$1,39 (te\'or. $\times$1,33) & $\times$1,44 (te\'or. $\times$1,25) & $\times$1,46 (te\'or. $\times$1,2) \\
\bottomrule
\end{tabular}}\par
\medskip

\end{table}

\Needspace{9\baselineskip}\paragraph{An\'alisis.}
\begin{itemize}
  \item \textbf{B\'usqueda lineal ($O(n)$).} Es la \'unica que crece como una recta de pendiente 1 en la gr\'afica
  log-log: el factor medido es $\times9$ a $\times13$ por salto frente al $\times10$ te\'orico. Por eso sirve
  como \emph{referencia}: no hace ning\'un supuesto sobre los datos, es la cota superior contra la que se mide
  cu\'anto ahorra aprovechar el orden. Con $n=1.000.000$ tarda 134~$\mu$s por b\'usqueda en Java (A) y 7,74~ms
  en Python (A); las dem\'as van de 70~ns a 635~ns en Java (A).
  \item \textbf{B\'usqueda por saltos ($O(\sqrt n)$).} Crece $\approx\times3$ por cada $\times10$ en $n$
  ($\times3{,}04$ y $\times3{,}19$ en Java, $\times3{,}37$ y $\times2{,}84$ en Python, en el \'ultimo salto),
  coincidiendo con el $\sqrt{10}=3{,}16$ te\'orico. Es mucho m\'as r\'apida que la lineal, pero queda entre
  entre 5 y 26 veces por detr\'as de las logar\'itmicas con $n=10^6$ (635~ns frente a 135~ns de la binaria en
  Java A; 33~$\mu$s frente a 1,26~$\mu$s en Python A).
  \item \textbf{Binaria, ternaria y exponencial ($O(\log n)$).} Casi no crecen: por cada $\times10$ en $n$
  el tiempo se multiplica por $\approx1{,}2$ a $1{,}8$ (con alg\'un valor aislado entre $\times0{,}8$ y
  $\times2{,}2$), cerca del $\times1{,}2$ a $\times1{,}33$ que predice $\log_2 n$ (de 10 a 20 iteraciones). Las desviaciones vienen de que son tiempos de decenas de nanosegundos,
  donde el ruido pesa, y de que al crecer el arreglo los \'ultimos accesos de cada b\'usqueda caen m\'as
  lejos en memoria. La ternaria hace $\log_3 n$ iteraciones pero con dos comparaciones cada una
  ($2\log_3 n\approx1{,}26\log_2 n$ comparaciones), as\'i que no gana a la binaria en comparaciones; en la
  pr\'actica quedaron empatadas dentro del ruido (en Python la ternaria es algo m\'as r\'apida porque ejecuta
  menos iteraciones del ciclo interpretado).
  \item \textbf{Interpolaci\'on ($O(\log\log n)$ promedio).} Fue la m\'as r\'apida con $n=10^6$ en los dos
  lenguajes y los dos equipos (70~ns y 110~ns en Java; 554~ns y 981~ns en Python), porque los datos son
  uniformes: estima la posici\'on y casi siempre acierta en pocos intentos. Su ventaja es modesta frente a la
  binaria (unas 2 veces) porque $\log\log n$ y $\log n$ ya son valores peque\~nos con $n=10^6$ ($\approx4{,}3$
  frente a $\approx20$); su factor de crecimiento ($\times1{,}0$ a $\times1{,}7$) no se distingue del de las
  logar\'itmicas dentro del ruido de medici\'on. Con datos no uniformes se degradar\'ia hacia $O(n)$.
  \item \textbf{A partir de qu\'e $n$ se nota la diferencia con la lineal.} En Python ya con $n=1.000$ la
  lineal es unas 12 veces m\'as lenta que la binaria (5,76~$\mu$s frente a 0,47~$\mu$s en A); en Java, con
  $n=1.000$ la diferencia es de solo 2,5 veces (116~ns frente a 47~ns), con $n=10.000$ ya son 11 veces y con
  $n=1.000.000$ son unas 1.000 veces (134~$\mu$s frente a 135~ns).
\end{itemize}

% ---------------------------------------------------------------------------- 2.2
\subsection{Compilado frente a interpretado}

\figpair{p2_compiled}{0.92}{Java frente a Python en b\'usqueda, $n=1.000.000$ (tiempo promedio por
b\'usqueda).}{fig:p2compiled}

\Needspace{9\baselineskip}\paragraph{An\'alisis.} El cociente Python/Java con $n=1.000.000$ es de $\times8$ a $\times10$ para las
b\'usquedas logar\'itmicas y de interpolaci\'on (binaria $\times9{,}3$, ternaria $\times8{,}4$, interpolaci\'on
$\times7{,}9$, exponencial $\times9{,}5$ en el equipo~A; $\times9{,}1$, $\times8{,}5$, $\times8{,}9$ y
$\times10{,}5$ en el B), pero de $\times52$ a $\times100$ para las que recorren muchos elementos: lineal
($\times57{,}7$ en A, $\times99{,}9$ en B) y por saltos ($\times52$ en A, $\times67{,}6$ en B). La raz\'on es la
misma del ordenamiento: cuando una b\'usqueda solo ejecuta $\approx20$ iteraciones, una parte importante del
tiempo en Java son esperas de memoria (cada sondeo cae lejos del anterior en un arreglo de 4~MB), un costo que
Python paga igual; solo se multiplica el costo de ejecutar las instrucciones, y la brecha se comprime. Cuando
hay $10^5$ o $10^6$ iteraciones consecutivas, la diferencia de costo \emph{por iteraci\'on} (c\'odigo de
m\'aquina frente a bytecode) se acumula casi por completo, y el ciclo lineal de Java, muy simple y sobre
memoria contigua, es adem\'as f\'acil de optimizar para el compilador. Es una interpretaci\'on de los datos; no
la medimos de forma aislada.

% ---------------------------------------------------------------------------- 2.3
\subsection{Posici\'on de la clave}

Las Tablas~\ref{tab:posJavaA}--\ref{tab:posPythonB} muestran el tiempo promedio por b\'usqueda cuando la misma
clave est\'a al inicio (\'indice 0), en el centro (\'indice $n/2$) y al final (\'indice $n-1$) del arreglo de
$n=1.000.000$, para cada equipo y lenguaje.

\begin{table}[htbp]
\centering
\caption{Equipo A, Java: tiempo promedio por b\'usqueda seg\'un la posici\'on de la clave ($n=1.000.000$).}\label{tab:posJavaA}
{\small
\begin{tabular}{lrrr}
\toprule
Algoritmo & Inicio & Centro & Final \\
\midrule
Lineal & 0,933 $\pm$ 0,4 ns & 72,6 $\pm$ 1,5 $\mu$s & 156 $\pm$ 11 $\mu$s \\
Binaria & 22,3 $\pm$ 1,1 ns & 21,2 $\pm$ 1,9 ns & 23,1 $\pm$ 0,59 ns \\
Ternaria & 23,6 $\pm$ 0,55 ns & 24,3 $\pm$ 1,3 ns & 33,7 $\pm$ 15 ns \\
Por saltos & 3,19 $\pm$ 0,18 ns & 300 $\pm$ 3,6 ns & 1,21 $\pm$ 0,058 $\mu$s \\
Interpolaci\'on & 3,12 $\pm$ 0,33 ns & 7,23 $\pm$ 0,2 ns & 3,28 $\pm$ 0,28 ns \\
Exponencial & 27,3 $\pm$ 20 ns & 48,3 $\pm$ 14 ns & 34,8 $\pm$ 4,2 ns \\
\bottomrule
\end{tabular}}\par

\end{table}
\begin{table}[htbp]
\centering
\caption{Equipo A, Python: tiempo promedio por b\'usqueda seg\'un la posici\'on de la clave ($n=1.000.000$).}\label{tab:posPythonA}
{\small
\begin{tabular}{lrrr}
\toprule
Algoritmo & Inicio & Centro & Final \\
\midrule
Lineal & 66,6 $\pm$ 2,5 ns & 5,12 $\pm$ 0,028 ms & 10,2 $\pm$ 0,061 ms \\
Binaria & 759 $\pm$ 10 ns & 793 $\pm$ 17 ns & 804 $\pm$ 8,7 ns \\
Ternaria & 696 $\pm$ 12 ns & 806 $\pm$ 13 ns & 770 $\pm$ 8,1 ns \\
Por saltos & 105 $\pm$ 1,3 ns & 18,9 $\pm$ 0,1 $\mu$s & 65 $\pm$ 1,9 $\mu$s \\
Interpolaci\'on & 108 $\pm$ 4,4 ns & 312 $\pm$ 3,4 ns & 125 $\pm$ 6,7 ns \\
Exponencial & 36,9 $\pm$ 2,1 ns & 1,03 $\pm$ 0,012 $\mu$s & 1,2 $\pm$ 0,0089 $\mu$s \\
\bottomrule
\end{tabular}}\par

\end{table}
\begin{table}[htbp]
\centering
\caption{Equipo B, Java: tiempo promedio por b\'usqueda seg\'un la posici\'on de la clave ($n=1.000.000$).}\label{tab:posJavaB}
{\small
\begin{tabular}{lrrr}
\toprule
Algoritmo & Inicio & Centro & Final \\
\midrule
Lineal & 0,864 $\pm$ 0,14 ns & 91 $\pm$ 3,4 $\mu$s & 174 $\pm$ 5 $\mu$s \\
Binaria & 33,3 $\pm$ 7,4 ns & 30,4 $\pm$ 0,099 ns & 39,8 $\pm$ 14 ns \\
Ternaria & 36 $\pm$ 2,3 ns & 44 $\pm$ 10 ns & 38,4 $\pm$ 2,7 ns \\
Por saltos & 4,89 $\pm$ 0,16 ns & 443 $\pm$ 4,3 ns & 1,57 $\pm$ 0,009 $\mu$s \\
Interpolaci\'on & 5,59 $\pm$ 0,28 ns & 30,3 $\pm$ 12 ns & 5,79 $\pm$ 0,06 ns \\
Exponencial & 41,9 $\pm$ 45 ns & 91,5 $\pm$ 36 ns & 60,6 $\pm$ 11 ns \\
\bottomrule
\end{tabular}}\par

\end{table}
\begin{table}[htbp]
\centering
\caption{Equipo B, Python: tiempo promedio por b\'usqueda seg\'un la posici\'on de la clave ($n=1.000.000$).}\label{tab:posPythonB}
{\small
\begin{tabular}{lrrr}
\toprule
Algoritmo & Inicio & Centro & Final \\
\midrule
Lineal & 95,8 $\pm$ 1,9 ns & 9,23 $\pm$ 0,21 ms & 18,7 $\pm$ 0,29 ms \\
Binaria & 1,31 $\pm$ 0,029 $\mu$s & 1,35 $\pm$ 0,02 $\mu$s & 1,4 $\pm$ 0,0084 $\mu$s \\
Ternaria & 1,2 $\pm$ 0,0093 $\mu$s & 1,45 $\pm$ 0,0091 $\mu$s & 1,4 $\pm$ 0,025 $\mu$s \\
Por saltos & 162 $\pm$ 5,3 ns & 29,5 $\pm$ 0,25 $\mu$s & 104 $\pm$ 2,6 $\mu$s \\
Interpolaci\'on & 147 $\pm$ 2,2 ns & 534 $\pm$ 6 ns & 186 $\pm$ 1,1 ns \\
Exponencial & 59,7 $\pm$ 1,2 ns & 1,66 $\pm$ 0,017 $\mu$s & 2 $\pm$ 0,0037 $\mu$s \\
\bottomrule
\end{tabular}}\par

\end{table}

\Needspace{9\baselineskip}\paragraph{An\'alisis.} El patr\'on es el mismo en los dos lenguajes y en los dos equipos:
\begin{itemize}
  \item \textbf{Cambia mucho en la lineal y en la de saltos.} Ambas recorren el arreglo desde el inicio, as\'i
  que su costo es proporcional a cu\'anto hay que avanzar. La lineal con la clave al inicio hace una sola
  comparaci\'on ($\approx1$~ns en Java, 67~ns en Python, que es b\'asicamente el costo de la llamada); en el
  centro revisa $n/2$ elementos (72,6~$\mu$s en Java A) y al final los $n$ (156~$\mu$s), de modo que el centro
  cuesta casi la mitad que el final. La de saltos con la clave al inicio la encuentra en el primer bloque
  (3~ns); en el centro da $\approx500$ saltos de $\sqrt n=1.000$ posiciones (300~ns), y al final da $\approx1.000$
  saltos y luego recorre un bloque completo (1,21~$\mu$s), por lo que final/centro es $\approx4$ y no 2.
  \item \textbf{Casi no cambia en la binaria y la ternaria.} Ambas siempre parten el intervalo, sin importar
  d\'onde est\'e la clave, y necesitan $\approx\log n$ pasos en cualquier posici\'on (22, 21 y 23~ns en Java A
  para la binaria). Un detalle: con $n=10^6$, el primer punto medio de la binaria es el \'indice $499.999$, y
  la clave <<centro>> est\'a en el \'indice $500.000$, as\'i que \emph{no} se encuentra en el primer intento y
  ese caso cuesta lo mismo que los otros dos; el mejor caso $O(1)$ de la binaria habr\'ia requerido ubicar la
  clave exactamente en el punto medio. Adem\'as, como se busca la misma clave 1.000 veces seguidas, el
  predictor de ramificaciones y la cach\'e se <<aprenden>> el recorrido, lo que hace estos tiempos algo
  menores que con claves variadas.
  \item \textbf{Interpolaci\'on: depende poco, pero no es constante.} Con la clave al inicio o al final la
  f\'ormula de interpolaci\'on cae exactamente en el extremo (3~ns); en el centro la estimaci\'on cae cerca pero
  no exactamente, por lo que son necesarias algunas iteraciones m\'as (7~ns en Java A; 312~ns frente a 108 y
  125~ns en Python A).
  \item \textbf{Exponencial: crece con el logaritmo de la posici\'on.} Si la clave est\'a en el \'indice 0 se
  detecta de inmediato; en el centro o al final debe duplicar el \'indice $\approx\log_2 i$ veces y luego hacer
  una b\'usqueda binaria en el \'ultimo intervalo, por eso el centro y el final cuestan casi lo mismo entre s\'i
  (1,03 y 1,20~$\mu$s en Python A). En Java, el valor <<inicio>> (27~ns en A, 42~ns en B) es mucho mayor que el
  $\approx1$~ns que se esperar\'ia de una salida inmediata; una hip\'otesis es que en el calentamiento el JIT
  nunca vio ejecutarse esa salida temprana y, la primera vez, tuvo que \emph{desoptimizar} y recompilar el
  m\'etodo, costo que al promediar sobre 1.000 b\'usquedas se nota. No lo comprobamos.
\end{itemize}

% ============================================================================== PUNTO 3
\section{Punto 3: an\'alisis comparativo}\label{sec:p3}

\subsection{Comparativos frente a no comparativos}

\paragraph{L\'imite inferior te\'orico.} Cualquier algoritmo de ordenamiento que solo obtiene informaci\'on del
orden relativo comparando pares de elementos puede modelarse como un \'arbol de decisi\'on binario: cada hoja
es una de las $n!$ permutaciones posibles de la entrada, y cada comparaci\'on es un nodo interno con dos
resultados. Un \'arbol binario con $n!$ hojas tiene altura al menos $\log_2(n!)$, y por la aproximaci\'on de
Stirling $\log_2(n!)=n\log_2 n-n\log_2 e+O(\log n)=\Theta(n\log n)$. Por tanto, \textbf{ning\'un ordenamiento por
comparaciones puede hacerlo en menos de $\Omega(n\log n)$ comparaciones en el peor caso}~\cite{clrs}. Con
$n=10^6$ eso significa al menos $\approx1{,}85\times10^7$ comparaciones.

\paragraph{Por qu\'e Radix Sort y Bucket Sort no est\'an sujetos a \'el.} El argumento anterior supone que la
\'unica operaci\'on que distingue elementos es la comparaci\'on. Radix Sort y Bucket Sort \emph{nunca comparan
dos elementos entre s\'i}: Radix extrae un d\'igito y lo usa como \'indice de un arreglo de conteo; Bucket
calcula la cubeta de un valor con una f\'ormula aritm\'etica sobre el propio valor. Es informaci\'on adicional
sobre los datos (que son enteros de rango acotado o de distribuci\'on conocida) y, a cambio, ordenan en tiempo
$O(d(n+b))$ y $O(n+k)$ respectivamente, es decir, lineal en $n$.

\begin{table}[htbp]
\centering
\caption{Tiempos de ordenamiento con $n=1.000.000$ (entrada aleatoria). Promedio $\pm$ desviaci\'on est\'andar.}\label{tab:q1}
\textbf{Equipo A}\par\smallskip
{\small
\begin{tabular}{llrr}
\toprule
Algoritmo & Tipo & Java & Python \\
\midrule
Insertion Sort & Comparativo & 36,7 $\pm$ 0,07 s & $\sim$\,1,91 h\,(est.) \\
Dual-Pivot QuickSort & Comparativo & 51,7 $\pm$ 0,35 ms & 623 $\pm$ 4,2 ms \\
Merge Sort & Comparativo & 77,7 $\pm$ 5,8 ms & 986 $\pm$ 14 ms \\
Heap Sort & Comparativo & 86 $\pm$ 1,1 ms & 1,97 $\pm$ 0,11 s \\
Radix Sort & No comparativo & 18,2 $\pm$ 0,57 ms & 1,34 $\pm$ 0,097 s \\
Bucket Sort & No comparativo & 9,08 $\pm$ 2,8 ms & 486 $\pm$ 50 ms \\
\bottomrule
\end{tabular}}\par
\medskip
\textbf{Equipo B}\par\smallskip
{\small
\begin{tabular}{llrr}
\toprule
Algoritmo & Tipo & Java & Python \\
\midrule
Insertion Sort & Comparativo & 55,3 $\pm$ 23 s & $\sim$\,3,81 h\,(est.) \\
Dual-Pivot QuickSort & Comparativo & 76,9 $\pm$ 2,4 ms & 1,05 $\pm$ 0,048 s \\
Merge Sort & Comparativo & 124 $\pm$ 6,1 ms & 1,75 $\pm$ 0,036 s \\
Heap Sort & Comparativo & 129 $\pm$ 4,1 ms & 3,06 $\pm$ 0,11 s \\
Radix Sort & No comparativo & 50,2 $\pm$ 1,4 ms & 1,89 $\pm$ 0,37 s \\
Bucket Sort & No comparativo & 32,5 $\pm$ 12 ms & 833 $\pm$ 74 ms \\
\bottomrule
\end{tabular}}\par

\end{table}

\paragraph{Comparaci\'on de tiempos.} En Java, Bucket Sort y Radix Sort fueron los dos m\'as r\'apidos en ambos
equipos: Bucket Sort tard\'o 9,08~ms (A) y 32,5~ms (B), unas 5,7 y 2,4 veces menos que el mejor de los
comparativos (Dual-Pivot QuickSort: 51,7~ms y 76,9~ms); Radix Sort, 18,2~ms y 50,2~ms (2,8 y 1,5 veces menos).
Merge Sort y Heap Sort, los m\'as lentos entre los $O(n\log n)$, tardaron entre 78~ms y 129~ms. En
Python la ventaja se reduce: Bucket Sort sigue siendo el mejor (0,486~s en A y 0,833~s en B), pero solo 1,3
veces m\'as r\'apido que Dual-Pivot QuickSort, y Radix Sort (1,34~s y 1,89~s) es \emph{m\'as lento} que Dual-Pivot
QuickSort y Merge Sort (aunque m\'as r\'apido que Heap Sort). La explicaci\'on es que Radix Sort hace
$d=8$ pasadas completas con tres ciclos cada una, y en un lenguaje interpretado el n\'umero de operaciones
elementales pesa m\'as que la ventaja asint\'otica ($8$ pasadas frente a $\log_2 10^6\approx20$ niveles): la
ventaja te\'orica de ser lineal solo se materializa si las operaciones son baratas, como en Java.

\paragraph{Supuesto de cada algoritmo y qu\'e pasar\'ia si no se cumple.}
\begin{itemize}
  \item \textbf{Radix Sort} supone claves enteras (o representables como cadenas de d\'igitos) con una
  cantidad \emph{fija y peque\~na} $d$ de d\'igitos; su costo es $O(d(n+b))$. En este taller $d=8$. Con n\'umeros
  de 16 d\'igitos se necesitar\'ian 16 pasadas, es decir, el tiempo \emph{se duplicar\'ia} en el mejor de los
  casos: a partir de nuestras mediciones, unos 36~ms en lugar de 18~ms en Java (A) y 1,34~s $\to\sim$2,7~s en
  Python (A), \textbf{estimaciones}, no mediciones. Adem\'as, 16 d\'igitos no caben en un \code{int} de Java (hay
  que usar \code{long}). Si $d$ creciera hasta $\sim\log n$ (por ejemplo, enteros de 64 bits), Radix Sort perder\'ia
  su ventaja sobre los comparativos.
  \item \textbf{Bucket Sort} supone que los valores est\'an distribuidos \emph{aproximadamente de manera
  uniforme} en el intervalo $[\min,\max]$, de modo que cada cubeta recibe $\approx1$ elemento. Si los valores
  se concentran (por ejemplo, casi todos cerca del m\'inimo y unos pocos valores extremos), la mayor\'ia cae en
  una o pocas cubetas, y el Insertion Sort interno trabaja sobre casi todo el arreglo: el costo pasa de $O(n)$ a
  $O(n^2)$. Con nuestras mediciones se puede estimar la magnitud: Insertion Sort sobre $10^6$ elementos aleatorios
  tarda 36,7~s en Java (A), frente a 9,08~ms del Bucket Sort con datos uniformes, es decir, alrededor de
  \textbf{4.000 veces m\'as} en el peor caso (estimaci\'on).
\end{itemize}

\subsection{\textquestiondown Ordenar o no ordenar?}\label{sec:q2}

Se tiene un arreglo \emph{desordenado} de 1.000.000 de elementos sobre el que se deben hacer $k$ b\'usquedas.
Hay dos estrategias: (A) $k$ b\'usquedas lineales, con costo $k\cdot t_{\text{lineal}}$; (B) ordenar una vez con el
mejor algoritmo medido y luego hacer $k$ b\'usquedas binarias, con costo $T_{\text{orden}}+k\cdot
t_{\text{binaria}}$. Ordenar conviene cuando
\begin{equation}
  T_{\text{orden}}+k\,t_{\text{binaria}}\;<\;k\,t_{\text{lineal}}
  \quad\Longleftrightarrow\quad
  k\;>\;k^*=\frac{T_{\text{orden}}}{t_{\text{lineal}}-t_{\text{binaria}}}.\label{eq:kstar}
\end{equation}
Como $T_{\text{orden}}$ se toma de ordenar un arreglo aleatorio de $10^6$ elementos, $t_{\text{lineal}}$ y
$t_{\text{binaria}}$ son los tiempos promedio por b\'usqueda medidos con $n=1.000.000$ (Punto~2). La b\'usqueda
lineal no aprovecha el orden, por lo que su costo por b\'usqueda es el mismo en el arreglo desordenado que en el
ordenado. El mejor algoritmo de ordenamiento medido fue, en los dos lenguajes y los dos equipos, \textbf{Bucket
Sort}.

\begin{table}[htbp]
\centering
\caption{C\'alculo del umbral $k^*$ con $n=1.000.000$.}\label{tab:q2}
{\footnotesize
\begin{tabular}{llllrrr}
\toprule
Equipo & Lenguaje & Mejor ordenamiento & $T_{\text{orden}}$ & $t_{\text{lineal}}$ & $t_{\text{binaria}}$ & $k^*$ \\
\midrule
A & Java & Bucket Sort & 9,08 ms & 134 $\mu$s & 135 ns & 67,71 $\Rightarrow$ 68 \\
A & Python & Bucket Sort & 486 ms & 7,74 ms & 1,26 $\mu$s & 62,83 $\Rightarrow$ 63 \\
B & Java & Bucket Sort & 32,5 ms & 144 $\mu$s & 229 ns & 226,1 $\Rightarrow$ 227 \\
B & Python & Bucket Sort & 833 ms & 14,4 ms & 2,08 $\mu$s & 57,85 $\Rightarrow$ 58 \\
\bottomrule
\end{tabular}}\par

\end{table}

\paragraph{C\'alculo (ecuaci\'on~\eqref{eq:kstar}).}
\begin{align*}
  \text{Java, A:}\quad & k^*=\frac{9{,}08\text{ ms}}{134{,}2\ \mu\text{s}-0{,}135\ \mu\text{s}}=\frac{9{,}08\times10^{-3}}{1{,}341\times10^{-4}}\approx67{,}7
    \;\Rightarrow\; k\ge68,\\
  \text{Python, A:}\quad & k^*=\frac{0{,}486\text{ s}}{7{,}74\text{ ms}-1{,}26\ \mu\text{s}}\approx62{,}8
    \;\Rightarrow\; k\ge63,\\
  \text{Java, B:}\quad & k^*=\frac{32{,}5\text{ ms}}{144{,}0\ \mu\text{s}-0{,}229\ \mu\text{s}}\approx226{,}1
    \;\Rightarrow\; k\ge227,\\
  \text{Python, B:}\quad & k^*=\frac{0{,}833\text{ s}}{14{,}4\text{ ms}-2{,}08\ \mu\text{s}}\approx57{,}9
    \;\Rightarrow\; k\ge58.
\end{align*}

\figpair{p3_q2}{0.82}{Costo total de cada estrategia seg\'un $k$ ($n=1.000.000$). Las l\'ineas verticales
punteadas marcan $k^*$.}{fig:q2}

\paragraph{Interpretaci\'on.} $k^*$ responde cu\'antas b\'usquedas hacen falta para que el costo \'unico de ordenar
se <<pague>>. Con menos de $k^*$ b\'usquedas (por ejemplo $k=1$, cuando ordenar costar\'ia 9~ms frente a 134~$\mu$s
de buscar linealmente en Java A), conviene la b\'usqueda lineal; con m\'as, conviene ordenar. A partir de $k^*$
la diferencia crece r\'apido: con $k=10^6$, la estrategia~A tarda $\approx134$~s en Java (A) y la B $\approx0{,}14$~s.
Los umbrales son del mismo orden de magnitud (decenas a pocos cientos) y por una raz\'on simple: ordenar con
Bucket Sort equivale, en tiempo, a unas 60--230 b\'usquedas lineales.
\begin{itemize}
  \item \textbf{Por qu\'e difiere entre lenguajes y equipos.} En Python, tanto ordenar como buscar linealmente son
  ciclos interpretados con costos similares por elemento, y el cociente $T_{\text{orden}}/t_{\text{lineal}}$ es
  estable (63 y 58 en los dos equipos). En Java, la b\'usqueda lineal es un ciclo simple sobre memoria contigua
  (muy r\'apido, $\approx0{,}2$~ns por elemento revisado) mientras que Bucket Sort hace escrituras dispersas; en el equipo~A
  ordenar equivale a 68 b\'usquedas lineales y en el B a 227, porque Bucket Sort empeor\'o mucho m\'as al pasar de A a
  B ($\times3{,}6$ en tiempo) que la b\'usqueda lineal ($\times1{,}07$): una pasada secuencial no depende tanto de
  la jerarqu\'ia de memoria como una escritura dispersa.
  \item \textbf{Salvedades.} (i) Los valores de $k^*$ usan el tiempo medio por b\'usqueda del conjunto de 1.000
  claves, mitad presentes y mitad ausentes; si las claves reales estuvieran casi siempre al inicio, la lineal
  ser\'ia mucho m\'as barata y $k^*$ mucho mayor. (ii) Si despu\'es de ordenar se necesitar\'an los datos en su orden
  original, habr\'ia que conservar una copia. (iii) Si los datos cambian entre b\'usquedas, hay que volver a ordenar
  y la ecuaci\'on~\eqref{eq:kstar} deja de aplicar.
\end{itemize}

% ============================================================================== COMPARACI\'ON ENTRE EQUIPOS
\section{Comparaci\'on entre los dos equipos}\label{sec:equipos}

La Tabla~\ref{tab:cross} muestra el cociente de tiempos equipo~B / equipo~A con $n=1.000.000$ (valores mayores
que $\times1$ indican que B fue m\'as lento).

\begin{table}[htbp]
\centering
\caption{Cociente de tiempos B/A con $n=1.000.000$ (entrada aleatoria para ordenamiento; 1.000 claves para
b\'usqueda).}\label{tab:cross}
{\small
\begin{tabular}{lrr}
\toprule
Algoritmo & Java & Python \\
\midrule
\multicolumn{3}{l}{\textit{Ordenamiento (aleatoria)}} \\
Insertion Sort & $\times$1,51 & $\sim$$\times$1,99 \\
Dual-Pivot QuickSort & $\times$1,49 & $\times$1,69 \\
Merge Sort & $\times$1,59 & $\times$1,77 \\
Heap Sort & $\times$1,5 & $\times$1,56 \\
Radix Sort & $\times$2,76 & $\times$1,41 \\
Bucket Sort & $\times$3,58 & $\times$1,71 \\
\multicolumn{3}{l}{\textit{B\'usqueda (1.000 claves)}} \\
Lineal & $\times$1,07 & $\times$1,86 \\
Binaria & $\times$1,69 & $\times$1,66 \\
Ternaria & $\times$1,83 & $\times$1,84 \\
Por saltos & $\times$1,32 & $\times$1,71 \\
Interpolaci\'on & $\times$1,57 & $\times$1,77 \\
Exponencial & $\times$1,64 & $\times$1,81 \\
\bottomrule
\end{tabular}}\par

\end{table}

\begin{itemize}
  \item \textbf{Mismas conclusiones, distinta escala.} En todos los casos el equipo~B fue m\'as lento (entre
  $\times1{,}1$ y $\times3{,}6$; la mayor\'ia entre $\times1{,}3$ y $\times2{,}0$), pero el orden de los
  algoritmos, los factores de crecimiento, los efectos del orden de entrada, la brecha Java--Python y los
  umbrales $k^*$ de Python son cualitativamente los mismos. Las conclusiones, por tanto, no dependen del
  equipo.
  \item \textbf{Donde m\'as se separan los equipos.} Los algoritmos cuyo costo est\'a dominado por accesos
  dispersos a memoria (Bucket Sort $\times3{,}6$ y Radix Sort $\times2{,}8$ en Java) se degradan m\'as en B que los
  dominados por c\'omputo (Merge, Heap, Dual-Pivot: $\times1{,}5$--$1{,}6$), y la b\'usqueda lineal en Java (acceso
  secuencial) casi no cambia ($\times1{,}07$). Es consistente con que B tenga una jerarqu\'ia de memoria menos
  favorable que A, aunque sin perfilar el hardware esa es solo una hip\'otesis.
  \item \textbf{Variabilidad.} Algunos tiempos de Java con $n=1.000$ (por ejemplo, Heap Sort en A) se alejan de
  lo esperado por efectos de calentamiento del JIT; con $n\ge10.000$ las desviaciones est\'andar son peque\~nas
  frente al promedio en la mayor\'ia de los casos (ver Anexo).
\end{itemize}

% ============================================================================== USO DE IA
\section{Uso de inteligencia artificial}

Se utiliz\'o \textbf{Claude Code} (Anthropic) como herramienta de apoyo en las siguientes partes del trabajo:
\begin{itemize}
  \item \textbf{Implementaci\'on de los algoritmos y de los programas de medici\'on:} primera versi\'on del c\'odigo
  en Java y Python, adaptada de los art\'iculos de GeeksforGeeks de la Tabla~\ref{tab:fuentes}, y de los
  programas de medici\'on (calentamiento, repeticiones, verificaci\'on y estimaci\'on de tiempos).
  \item \textbf{Generaci\'on de gr\'aficas y tablas:} scripts en Python (matplotlib) que leen los CSV y
  producen las gr\'aficas; script que genera las tablas en \LaTeX{} de este documento a partir de los CSV de los
  dos equipos.
  \item \textbf{Investigaci\'on y redacci\'on:} consulta de las complejidades te\'oricas y redacci\'on de un
  borrador de los an\'alisis y del documento \LaTeX{}, a partir de los resultados medidos.
\end{itemize}
Las mediciones se ejecutaron en los dos equipos por el grupo; la herramienta no gener\'o ni modific\'o ning\'un
dato experimental. \pendiente{El grupo debe revisar este texto y ajustarlo para que describa con exactitud
el uso real que hizo, y revisar con sus propias palabras los an\'alisis antes de entregar.}

% ============================================================================== CONCLUSIONES
\section{Conclusiones}

\begin{enumerate}
  \item \textbf{El comportamiento coincide con la teor\'ia.} Al multiplicar $n$ por 10, Insertion Sort multiplic\'o
  su tiempo por $\approx100$ ($O(n^2)$), Merge Sort, Heap Sort y Dual-Pivot QuickSort por $\approx11$--$13$
  ($O(n\log n)$) y Radix Sort por $\approx10$ ($O(n)$), en Java y en Python y en ambos equipos. Las
  desviaciones (hasta $\times17$--$19$ en Radix y Bucket Sort con $n=10^6$) se explican por la jerarqu\'ia de
  memoria y por los costos del JIT con $n$ peque\~no, no por un error de complejidad.
  \item \textbf{El orden de entrada puede cambiar la complejidad.} Insertion Sort pas\'o de 361~ms a 53~$\mu$s
  ($n=10^5$, Java, A) con entrada ascendente, pero fue el doble de lento con descendente. Dual-Pivot QuickSort
  con pivotes en los extremos degener\'o a $O(n^2)$ con datos ordenados y fall\'o por desbordamiento de pila
  (Java, $n=10^5$) o por el l\'imite de recursi\'on (Python, $n=10^4$). Merge, Heap, Radix y Bucket Sort no
  cambiaron de complejidad con el orden.
  \item \textbf{Mejores algoritmos.} Con $n=10^6$ aleatorio, Bucket Sort fue el m\'as r\'apido (9,08~ms en Java A;
  0,486~s en Python A), seguido de Radix Sort en Java y de Dual-Pivot QuickSort en Python. Entre los comparativos,
  Dual-Pivot QuickSort gan\'o con datos aleatorios, pero es el \'unico que no es seguro con datos ordenados. Los no
  comparativos superan al l\'imite $\Omega(n\log n)$ porque no comparan, a cambio de supuestos sobre los datos
  (d\'igitos fijos, distribuci\'on uniforme).
  \item \textbf{Compilado frente a interpretado.} Java fue entre $\times8$ y $\times250$ m\'as r\'apido. La
  brecha es peque\~na ($\times8$--$14$) en b\'usquedas logar\'itmicas y en Merge Sort/Dual-Pivot QuickSort, y
  grande ($\times50$--$250$) en lo que repite muchas iteraciones simples (b\'usqueda lineal y por saltos,
  Insertion, Radix y Bucket Sort).
  \item \textbf{B\'usqueda.} Con datos uniformes, la interpolaci\'on fue la m\'as r\'apida; binaria, ternaria y
  exponencial quedaron empatadas dentro del ruido; la de saltos creci\'o como $\sqrt n$ y la lineal, como $n$. La
  posici\'on de la clave solo importa para los m\'etodos que recorren desde el inicio (lineal, saltos y, en
  menor medida, exponencial e interpolaci\'on).
  \item \textbf{\textquestiondown Ordenar o no ordenar?} Para $n=10^6$ conviene ordenar con Bucket Sort y usar
  b\'usqueda binaria a partir de $k^*\approx58$--$227$ b\'usquedas, seg\'un lenguaje y equipo.
  \item \textbf{Limitaciones.} Cinco repeticiones por medici\'on; los valores estimados (Insertion Sort en Python
  con $n=10^6$ y Dual-Pivot QuickSort con entrada ordenada) escalan con un modelo te\'orico y deben leerse como
  aproximaciones; las explicaciones sobre cach\'e, JIT y predicci\'on de ramificaciones son hip\'otesis razonables
  que no verificamos con herramientas de perfilado.
\end{enumerate}

% ============================================================================== BIBLIOGRAF\'IA
\begin{thebibliography}{9}
\bibitem{clrs} T.~H.~Cormen, C.~E.~Leiserson, R.~L.~Rivest y C.~Stein, \emph{Introduction to Algorithms},
  3.\textsuperscript{a} ed., MIT Press, 2009. Cap.~8, Teorema~8.1 (cota inferior de ordenamiento por comparaciones).
\bibitem{gfg} GeeksforGeeks, art\'iculos de algoritmos de ordenamiento y b\'usqueda (ver Tabla~\ref{tab:fuentes}).
  \url{https://www.geeksforgeeks.org/}
\end{thebibliography}

% ================================================================================= ANEXO
\clearpage
\appendix
\section{Tablas completas de resultados}\label{anexo}

Promedio $\pm$ desviaci\'on est\'andar muestral de 5 repeticiones. <<$\sim$\,\ldots\,(est.)>> indica un valor
estimado, no medido.

\subsection{Equipo A}

\begin{table}[H]
\centering
\caption{Equipo A, ordenamiento en Java.}{\scriptsize
\begin{tabular}{llrrrr}
\toprule
Algoritmo & Entrada & $n=1.000$ & $n=10.000$ & $n=100.000$ & $n=1.000.000$ \\
\midrule
Insertion Sort & Aleatoria & 47,1 $\pm$ 7 $\mu$s & 3,5 $\pm$ 0,062 ms & 361 $\pm$ 0,75 ms & 36,7 $\pm$ 0,07 s \\
 & Ascendente & 3,18 $\pm$ 2,7 $\mu$s & 14,1 $\pm$ 11 $\mu$s & 53,1 $\pm$ 0,59 $\mu$s & 572 $\pm$ 77 $\mu$s \\
 & Descendente & 91,9 $\pm$ 2,8 $\mu$s & 7,47 $\pm$ 0,12 ms & 729 $\pm$ 7,3 ms & 1,23 $\pm$ 0,00037 min \\
Dual-Pivot QuickSort & Aleatoria & 44,6 $\pm$ 13 $\mu$s & 417 $\pm$ 70 $\mu$s & 4,67 $\pm$ 0,26 ms & 51,7 $\pm$ 0,35 ms \\
 & Ascendente & 138 $\pm$ 3,8 $\mu$s & 14,6 $\pm$ 0,12 ms & $\sim$\,1,46 s\,(est.) & $\sim$\,2,43 min\,(est.) \\
 & Descendente & 135 $\pm$ 1,6 $\mu$s & 15,1 $\pm$ 0,44 ms & $\sim$\,1,51 s\,(est.) & $\sim$\,2,52 min\,(est.) \\
Merge Sort & Aleatoria & 46,8 $\pm$ 7,6 $\mu$s & 537 $\pm$ 19 $\mu$s & 6,47 $\pm$ 0,17 ms & 77,7 $\pm$ 5,8 ms \\
 & Ascendente & 20,9 $\pm$ 1,6 $\mu$s & 174 $\pm$ 4,3 $\mu$s & 1,9 $\pm$ 0,049 ms & 28,8 $\pm$ 16 ms \\
 & Descendente & 21,1 $\pm$ 2,5 $\mu$s & 176 $\pm$ 3,1 $\mu$s & 1,89 $\pm$ 0,037 ms & 21 $\pm$ 0,19 ms \\
Heap Sort & Aleatoria & 210 $\pm$ 364,17 $\mu$s & 572 $\pm$ 40 $\mu$s & 6,81 $\pm$ 0,12 ms & 86 $\pm$ 1,1 ms \\
 & Ascendente & 34,9 $\pm$ 2,1 $\mu$s & 372 $\pm$ 8,5 $\mu$s & 4,01 $\pm$ 0,13 ms & 46,2 $\pm$ 3,2 ms \\
 & Descendente & 32,2 $\pm$ 2,9 $\mu$s & 387 $\pm$ 9 $\mu$s & 4,19 $\pm$ 0,079 ms & 45,4 $\pm$ 0,16 ms \\
Radix Sort & Aleatoria & 26,2 $\pm$ 4,4 $\mu$s & 200 $\pm$ 5,4 $\mu$s & 1,93 $\pm$ 0,1 ms & 18,2 $\pm$ 0,57 ms \\
 & Ascendente & 22,9 $\pm$ 2,3 $\mu$s & 188 $\pm$ 6,6 $\mu$s & 2 $\pm$ 0,17 ms & 18,4 $\pm$ 0,18 ms \\
 & Descendente & 22,8 $\pm$ 1,5 $\mu$s & 201 $\pm$ 10 $\mu$s & 2,49 $\pm$ 0,42 ms & 17,9 $\pm$ 0,22 ms \\
Bucket Sort & Aleatoria & 26,3 $\pm$ 11 $\mu$s & 167 $\pm$ 64 $\mu$s & 1,96 $\pm$ 0,88 ms & 9,08 $\pm$ 2,8 ms \\
 & Ascendente & 10,8 $\pm$ 6,2 $\mu$s & 52,6 $\pm$ 5,1 $\mu$s & 729 $\pm$ 303,30 $\mu$s & 5,8 $\pm$ 0,91 ms \\
 & Descendente & 8,97 $\pm$ 1,4 $\mu$s & 58,8 $\pm$ 6,9 $\mu$s & 577 $\pm$ 26 $\mu$s & 6,04 $\pm$ 0,12 ms \\
\bottomrule
\end{tabular}}\par

\end{table}
\begin{table}[H]
\centering
\caption{Equipo A, ordenamiento en Python.}{\scriptsize
\begin{tabular}{llrrrr}
\toprule
Algoritmo & Entrada & $n=1.000$ & $n=10.000$ & $n=100.000$ & $n=1.000.000$ \\
\midrule
Insertion Sort & Aleatoria & 4,95 $\pm$ 0,033 ms & 611 $\pm$ 12 ms & 1,15 $\pm$ 0,076 min & $\sim$\,1,91 h\,(est.) \\
 & Ascendente & 30,9 $\pm$ 1 $\mu$s & 300 $\pm$ 4,1 $\mu$s & 3,22 $\pm$ 0,21 ms & 30 $\pm$ 0,16 ms \\
 & Descendente & 9,82 $\pm$ 0,17 ms & 1,16 $\pm$ 0,013 s & 2,07 $\pm$ 0,058 min & $\sim$\,3,45 h\,(est.) \\
Dual-Pivot QuickSort & Aleatoria & 263 $\pm$ 18 $\mu$s & 3,91 $\pm$ 0,051 ms & 47,5 $\pm$ 0,15 ms & 623 $\pm$ 4,2 ms \\
 & Ascendente & 4,48 $\pm$ 0,032 ms & $\sim$\,448 ms\,(est.) & $\sim$\,44,8 s\,(est.) & $\sim$\,1,24 h\,(est.) \\
 & Descendente & 4,58 $\pm$ 0,032 ms & $\sim$\,458 ms\,(est.) & $\sim$\,45,8 s\,(est.) & $\sim$\,1,27 h\,(est.) \\
Merge Sort & Aleatoria & 472 $\pm$ 16 $\mu$s & 6,36 $\pm$ 0,038 ms & 77,6 $\pm$ 0,27 ms & 986 $\pm$ 14 ms \\
 & Ascendente & 367 $\pm$ 3 $\mu$s & 4,97 $\pm$ 0,022 ms & 58,8 $\pm$ 0,068 ms & 678 $\pm$ 6,1 ms \\
 & Descendente & 360 $\pm$ 3 $\mu$s & 4,87 $\pm$ 0,019 ms & 57,9 $\pm$ 0,18 ms & 665 $\pm$ 2,1 ms \\
Heap Sort & Aleatoria & 658 $\pm$ 15 $\mu$s & 9,81 $\pm$ 0,15 ms & 126 $\pm$ 0,62 ms & 1,97 $\pm$ 0,11 s \\
 & Ascendente & 670 $\pm$ 10 $\mu$s & 9,45 $\pm$ 0,082 ms & 121 $\pm$ 0,91 ms & 1,46 $\pm$ 0,012 s \\
 & Descendente & 589 $\pm$ 6,8 $\mu$s & 8,52 $\pm$ 0,03 ms & 111 $\pm$ 0,45 ms & 1,41 $\pm$ 0,029 s \\
Radix Sort & Aleatoria & 655 $\pm$ 102,37 $\mu$s & 7,23 $\pm$ 0,49 ms & 69,6 $\pm$ 0,51 ms & 1,34 $\pm$ 0,097 s \\
 & Ascendente & 635 $\pm$ 28 $\mu$s & 7 $\pm$ 0,1 ms & 68,3 $\pm$ 0,96 ms & 924 $\pm$ 30 ms \\
 & Descendente & 611 $\pm$ 4,7 $\mu$s & 6,88 $\pm$ 0,064 ms & 68,3 $\pm$ 1,1 ms & 1,07 $\pm$ 0,13 s \\
Bucket Sort & Aleatoria & 186 $\pm$ 13 $\mu$s & 2,38 $\pm$ 0,063 ms & 26,3 $\pm$ 1,2 ms & 486 $\pm$ 50 ms \\
 & Ascendente & 174 $\pm$ 8,2 $\mu$s & 2,17 $\pm$ 0,18 ms & 21,3 $\pm$ 0,24 ms & 224 $\pm$ 6,6 ms \\
 & Descendente & 183 $\pm$ 2,1 $\mu$s & 2,18 $\pm$ 0,023 ms & 22,2 $\pm$ 0,2 ms & 262 $\pm$ 36 ms \\
\bottomrule
\end{tabular}}\par

\end{table}
\begin{table}[H]
\centering
\caption{Equipo A, b\'usqueda en Java (tiempo por b\'usqueda, 1.000 claves).}{\scriptsize
\begin{tabular}{lrrrr}
\toprule
Algoritmo & $n=1.000$ & $n=10.000$ & $n=100.000$ & $n=1.000.000$ \\
\midrule
Lineal & 116 $\pm$ 3,3 ns & 1,11 $\pm$ 0,028 $\mu$s & 11,1 $\pm$ 0,32 $\mu$s & 134 $\pm$ 20 $\mu$s \\
Binaria & 47,1 $\pm$ 6,1 ns & 102 $\pm$ 90 ns & 84,7 $\pm$ 13 ns & 135 $\pm$ 50 ns \\
Ternaria & 37,3 $\pm$ 12 ns & 54,3 $\pm$ 11 ns & 74 $\pm$ 18 ns & 126 $\pm$ 54 ns \\
Por saltos & 109 $\pm$ 59 ns & 81,4 $\pm$ 4,6 ns & 209 $\pm$ 11 ns & 635 $\pm$ 46 ns \\
Interpolaci\'on & 28,9 $\pm$ 2,2 ns & 41,1 $\pm$ 9,6 ns & 41,8 $\pm$ 15 ns & 70 $\pm$ 42 ns \\
Exponencial & 45,6 $\pm$ 7,7 ns & 51 $\pm$ 8,2 ns & 101 $\pm$ 14 ns & 156 $\pm$ 52 ns \\
\bottomrule
\end{tabular}}\par

\end{table}
\begin{table}[H]
\centering
\caption{Equipo A, b\'usqueda en Python (tiempo por b\'usqueda, 1.000 claves).}{\scriptsize
\begin{tabular}{lrrrr}
\toprule
Algoritmo & $n=1.000$ & $n=10.000$ & $n=100.000$ & $n=1.000.000$ \\
\midrule
Lineal & 5,76 $\pm$ 0,23 $\mu$s & 77,2 $\pm$ 2,6 $\mu$s & 779 $\pm$ 8,3 $\mu$s & 7,74 $\pm$ 0,099 ms \\
Binaria & 471 $\pm$ 22 ns & 696 $\pm$ 52 ns & 954 $\pm$ 198,41 ns & 1,26 $\pm$ 0,39 $\mu$s \\
Ternaria & 446 $\pm$ 17 ns & 621 $\pm$ 24 ns & 832 $\pm$ 93 ns & 1,06 $\pm$ 0,26 $\mu$s \\
Por saltos & 905 $\pm$ 12 ns & 3,15 $\pm$ 0,017 $\mu$s & 9,8 $\pm$ 0,072 $\mu$s & 33 $\pm$ 0,74 $\mu$s \\
Interpolaci\'on & 328 $\pm$ 9,8 ns & 421 $\pm$ 18 ns & 500 $\pm$ 96 ns & 554 $\pm$ 131,21 ns \\
Exponencial & 577 $\pm$ 13 ns & 831 $\pm$ 23 ns & 1,16 $\pm$ 0,13 $\mu$s & 1,49 $\pm$ 0,28 $\mu$s \\
\bottomrule
\end{tabular}}\par

\end{table}
\begin{table}[H]
\centering
\caption{Equipo A: registro de valores estimados y fallos.}\label{tab:issuesA}{\tiny
\begin{tabular}{lllrlp{4.6cm}}
\toprule
Lenguaje & Algoritmo & Entrada & $n$ & Estado & Nota \\
\midrule
Java & DualPivotQuickSort & asc & 100.000 & ERROR\_StackOverflowError & StackOverflowError at n=100000; scaled from n=10000 \\
Java & DualPivotQuickSort & asc & 1.000.000 & ESTIMATED & same failure expected (recursion depth grows with n); scaled from n=10000 \\
Java & DualPivotQuickSort & desc & 100.000 & ERROR\_StackOverflowError & StackOverflowError at n=100000; scaled from n=10000 \\
Java & DualPivotQuickSort & desc & 1.000.000 & ESTIMATED & same failure expected (recursion depth grows with n); scaled from n=10000 \\
Python & DualPivotQuickSort & asc & 10.000 & ERROR\_RecursionError & RecursionError at n=10000 (limit 1000); scaled from n=1000 \\
Python & DualPivotQuickSort & asc & 100.000 & ESTIMATED & same failure expected (recursion depth grows with n); scaled from n=1000 \\
Python & DualPivotQuickSort & asc & 1.000.000 & ESTIMATED & same failure expected (recursion depth grows with n); scaled from n=1000 \\
Python & DualPivotQuickSort & desc & 10.000 & ERROR\_RecursionError & RecursionError at n=10000 (limit 1000); scaled from n=1000 \\
Python & DualPivotQuickSort & desc & 100.000 & ESTIMATED & same failure expected (recursion depth grows with n); scaled from n=1000 \\
Python & DualPivotQuickSort & desc & 1.000.000 & ESTIMATED & same failure expected (recursion depth grows with n); scaled from n=1000 \\
Python & InsertionSort & desc & 1.000.000 & ESTIMATED & predicted 12413s $>$ limit 300s; scaled from n=100000 \\
Python & InsertionSort & random & 1.000.000 & ESTIMATED & predicted 6890s $>$ limit 300s; scaled from n=100000 \\
\bottomrule
\end{tabular}}\par

\end{table}

\subsection{Equipo B}

\begin{table}[H]
\centering
\caption{Equipo B, ordenamiento en Java.}{\scriptsize
\begin{tabular}{llrrrr}
\toprule
Algoritmo & Entrada & $n=1.000$ & $n=10.000$ & $n=100.000$ & $n=1.000.000$ \\
\midrule
Insertion Sort & Aleatoria & 49,4 $\pm$ 2,4 $\mu$s & 3,47 $\pm$ 0,24 ms & 372 $\pm$ 1,9 ms & 55,3 $\pm$ 23 s \\
 & Ascendente & 3,4 $\pm$ 0,78 $\mu$s & 8,67 $\pm$ 0,26 $\mu$s & 76,6 $\pm$ 10 $\mu$s & 813 $\pm$ 84 $\mu$s \\
 & Descendente & 87,1 $\pm$ 7,3 $\mu$s & 7,19 $\pm$ 0,094 ms & 792 $\pm$ 18 ms & 1,55 $\pm$ 0,02 min \\
Dual-Pivot QuickSort & Aleatoria & 51,1 $\pm$ 4,8 $\mu$s & 598 $\pm$ 73 $\mu$s & 6,61 $\pm$ 0,023 ms & 76,9 $\pm$ 2,4 ms \\
 & Ascendente & 178 $\pm$ 8,4 $\mu$s & 16,4 $\pm$ 0,054 ms & $\sim$\,1,64 s\,(est.) & $\sim$\,2,73 min\,(est.) \\
 & Descendente & 177 $\pm$ 0,89 $\mu$s & 17,4 $\pm$ 0,22 ms & $\sim$\,1,74 s\,(est.) & $\sim$\,2,9 min\,(est.) \\
Merge Sort & Aleatoria & 75 $\pm$ 4,3 $\mu$s & 810 $\pm$ 25 $\mu$s & 9,94 $\pm$ 0,16 ms & 124 $\pm$ 6,1 ms \\
 & Ascendente & 39,9 $\pm$ 17 $\mu$s & 277 $\pm$ 14 $\mu$s & 3,2 $\pm$ 0,037 ms & 42,2 $\pm$ 2,3 ms \\
 & Descendente & 38,2 $\pm$ 1,5 $\mu$s & 321 $\pm$ 6,1 $\mu$s & 3,8 $\pm$ 0,058 ms & 48,3 $\pm$ 2,4 ms \\
Heap Sort & Aleatoria & 84,5 $\pm$ 31 $\mu$s & 928 $\pm$ 24 $\mu$s & 10,6 $\pm$ 0,065 ms & 129 $\pm$ 4,1 ms \\
 & Ascendente & 41,5 $\pm$ 3,3 $\mu$s & 597 $\pm$ 91 $\mu$s & 5,65 $\pm$ 0,038 ms & 62,3 $\pm$ 1,2 ms \\
 & Descendente & 49,4 $\pm$ 3,4 $\mu$s & 535 $\pm$ 6,2 $\mu$s & 5,81 $\pm$ 0,01 ms & 64,1 $\pm$ 1,2 ms \\
Radix Sort & Aleatoria & 74,2 $\pm$ 31 $\mu$s & 481 $\pm$ 26 $\mu$s & 4,59 $\pm$ 0,046 ms & 50,2 $\pm$ 1,4 ms \\
 & Ascendente & 49,8 $\pm$ 2 $\mu$s & 451 $\pm$ 4,7 $\mu$s & 4,59 $\pm$ 0,04 ms & 50,5 $\pm$ 1,7 ms \\
 & Descendente & 48 $\pm$ 1,4 $\mu$s & 448 $\pm$ 20 $\mu$s & 4,51 $\pm$ 0,035 ms & 48,3 $\pm$ 0,7 ms \\
Bucket Sort & Aleatoria & 45,8 $\pm$ 6,3 $\mu$s & 224 $\pm$ 47 $\mu$s & 1,87 $\pm$ 0,17 ms & 32,5 $\pm$ 12 ms \\
 & Ascendente & 14,7 $\pm$ 2,6 $\mu$s & 111 $\pm$ 20 $\mu$s & 1,01 $\pm$ 0,071 ms & 12,5 $\pm$ 0,61 ms \\
 & Descendente & 12,9 $\pm$ 2,4 $\mu$s & 102 $\pm$ 3,1 $\mu$s & 1,03 $\pm$ 0,02 ms & 13,9 $\pm$ 1,5 ms \\
\bottomrule
\end{tabular}}\par

\end{table}
\begin{table}[H]
\centering
\caption{Equipo B, ordenamiento en Python.}{\scriptsize
\begin{tabular}{llrrrr}
\toprule
Algoritmo & Entrada & $n=1.000$ & $n=10.000$ & $n=100.000$ & $n=1.000.000$ \\
\midrule
Insertion Sort & Aleatoria & 9,74 $\pm$ 0,1 ms & 1,01 $\pm$ 0,013 s & 2,29 $\pm$ 0,45 min & $\sim$\,3,81 h\,(est.) \\
 & Ascendente & 51,9 $\pm$ 8,4 $\mu$s & 512 $\pm$ 7,1 $\mu$s & 5,4 $\pm$ 0,57 ms & 53,4 $\pm$ 0,3 ms \\
 & Descendente & 18,1 $\pm$ 0,062 ms & 1,98 $\pm$ 0,049 s & 3,69 $\pm$ 0,39 min & $\sim$\,6,15 h\,(est.) \\
Dual-Pivot QuickSort & Aleatoria & 434 $\pm$ 11 $\mu$s & 6,19 $\pm$ 0,036 ms & 77,3 $\pm$ 1,2 ms & 1,05 $\pm$ 0,048 s \\
 & Ascendente & 9,11 $\pm$ 0,36 ms & $\sim$\,911 ms\,(est.) & $\sim$\,1,52 min\,(est.) & $\sim$\,2,53 h\,(est.) \\
 & Descendente & 8,39 $\pm$ 0,074 ms & $\sim$\,839 ms\,(est.) & $\sim$\,1,4 min\,(est.) & $\sim$\,2,33 h\,(est.) \\
Merge Sort & Aleatoria & 802 $\pm$ 22 $\mu$s & 10,4 $\pm$ 0,11 ms & 130 $\pm$ 2,6 ms & 1,75 $\pm$ 0,036 s \\
 & Ascendente & 718 $\pm$ 35 $\mu$s & 9,21 $\pm$ 0,13 ms & 110 $\pm$ 2,5 ms & 1,24 $\pm$ 0,044 s \\
 & Descendente & 687 $\pm$ 23 $\mu$s & 8,9 $\pm$ 0,12 ms & 107 $\pm$ 3,4 ms & 1,2 $\pm$ 0,03 s \\
Heap Sort & Aleatoria & 1,17 $\pm$ 0,024 ms & 16,9 $\pm$ 0,22 ms & 223 $\pm$ 4,9 ms & 3,06 $\pm$ 0,11 s \\
 & Ascendente & 1,08 $\pm$ 0,014 ms & 15,6 $\pm$ 0,27 ms & 206 $\pm$ 2,5 ms & 2,48 $\pm$ 0,017 s \\
 & Descendente & 915 $\pm$ 28 $\mu$s & 14,4 $\pm$ 0,24 ms & 188 $\pm$ 2,3 ms & 2,34 $\pm$ 0,014 s \\
Radix Sort & Aleatoria & 792 $\pm$ 52 $\mu$s & 9,89 $\pm$ 0,65 ms & 109 $\pm$ 2,3 ms & 1,89 $\pm$ 0,37 s \\
 & Ascendente & 1,49 $\pm$ 0,11 ms & 15 $\pm$ 0,26 ms & 166 $\pm$ 8,2 ms & 1,86 $\pm$ 0,16 s \\
 & Descendente & 1,36 $\pm$ 0,023 ms & 13,8 $\pm$ 0,13 ms & 151 $\pm$ 7,4 ms & 1,7 $\pm$ 0,11 s \\
Bucket Sort & Aleatoria & 450 $\pm$ 27 $\mu$s & 4,59 $\pm$ 0,014 ms & 57 $\pm$ 1,3 ms & 833 $\pm$ 74 ms \\
 & Ascendente & 578 $\pm$ 404,55 $\mu$s & 4,17 $\pm$ 0,019 ms & 44,4 $\pm$ 0,27 ms & 473 $\pm$ 16 ms \\
 & Descendente & 411 $\pm$ 35 $\mu$s & 4,29 $\pm$ 0,079 ms & 45,7 $\pm$ 0,45 ms & 481 $\pm$ 12 ms \\
\bottomrule
\end{tabular}}\par

\end{table}
\begin{table}[H]
\centering
\caption{Equipo B, b\'usqueda en Java (tiempo por b\'usqueda, 1.000 claves).}{\scriptsize
\begin{tabular}{lrrrr}
\toprule
Algoritmo & $n=1.000$ & $n=10.000$ & $n=100.000$ & $n=1.000.000$ \\
\midrule
Lineal & 150 $\pm$ 3,7 ns & 1,68 $\pm$ 0,38 $\mu$s & 15,3 $\pm$ 1,6 $\mu$s & 144 $\pm$ 7,1 $\mu$s \\
Binaria & 67,6 $\pm$ 2,5 ns & 90,7 $\pm$ 1,6 ns & 135 $\pm$ 13 ns & 229 $\pm$ 43 ns \\
Ternaria & 75,7 $\pm$ 11 ns & 104 $\pm$ 12 ns & 127 $\pm$ 11 ns & 231 $\pm$ 58 ns \\
Por saltos & 215 $\pm$ 126,90 ns & 110 $\pm$ 6 ns & 262 $\pm$ 28 ns & 837 $\pm$ 34 ns \\
Interpolaci\'on & 50,8 $\pm$ 3,7 ns & 60,1 $\pm$ 4,3 ns & 86,6 $\pm$ 11 ns & 110 $\pm$ 24 ns \\
Exponencial & 72,7 $\pm$ 1,3 ns & 96,7 $\pm$ 2,1 ns & 151 $\pm$ 29 ns & 255 $\pm$ 36 ns \\
\bottomrule
\end{tabular}}\par

\end{table}
\begin{table}[H]
\centering
\caption{Equipo B, b\'usqueda en Python (tiempo por b\'usqueda, 1.000 claves).}{\scriptsize
\begin{tabular}{lrrrr}
\toprule
Algoritmo & $n=1.000$ & $n=10.000$ & $n=100.000$ & $n=1.000.000$ \\
\midrule
Lineal & 14,4 $\pm$ 0,035 $\mu$s & 155 $\pm$ 3,3 $\mu$s & 1,48 $\pm$ 0,073 ms & 14,4 $\pm$ 0,45 ms \\
Binaria & 749 $\pm$ 15 ns & 1,08 $\pm$ 0,026 $\mu$s & 1,47 $\pm$ 0,17 $\mu$s & 2,08 $\pm$ 0,32 $\mu$s \\
Ternaria & 736 $\pm$ 17 ns & 1,04 $\pm$ 0,018 $\mu$s & 1,42 $\pm$ 0,14 $\mu$s & 1,95 $\pm$ 0,31 $\mu$s \\
Por saltos & 1,98 $\pm$ 0,21 $\mu$s & 6,37 $\pm$ 0,2 $\mu$s & 19,9 $\pm$ 0,46 $\mu$s & 56,6 $\pm$ 2,5 $\mu$s \\
Interpolaci\'on & 553 $\pm$ 7,9 ns & 677 $\pm$ 29 ns & 808 $\pm$ 144,14 ns & 981 $\pm$ 197,13 ns \\
Exponencial & 926 $\pm$ 26 ns & 1,28 $\pm$ 0,021 $\mu$s & 1,84 $\pm$ 0,22 $\mu$s & 2,69 $\pm$ 0,34 $\mu$s \\
\bottomrule
\end{tabular}}\par

\end{table}
\begin{table}[H]
\centering
\caption{Equipo B: registro de valores estimados y fallos.}\label{tab:issuesB}{\tiny
\begin{tabular}{lllrlp{4.6cm}}
\toprule
Lenguaje & Algoritmo & Entrada & $n$ & Estado & Nota \\
\midrule
Java & DualPivotQuickSort & asc & 100.000 & ERROR\_StackOverflowError & StackOverflowError at n=100000; scaled from n=10000 \\
Java & DualPivotQuickSort & asc & 1.000.000 & ESTIMATED & same failure expected (recursion depth grows with n); scaled from n=10000 \\
Java & DualPivotQuickSort & desc & 100.000 & ERROR\_StackOverflowError & StackOverflowError at n=100000; scaled from n=10000 \\
Java & DualPivotQuickSort & desc & 1.000.000 & ESTIMATED & same failure expected (recursion depth grows with n); scaled from n=10000 \\
Python & DualPivotQuickSort & asc & 10.000 & ERROR\_RecursionError & RecursionError at n=10000 (limit 1000); scaled from n=1000 \\
Python & DualPivotQuickSort & asc & 100.000 & ESTIMATED & same failure expected (recursion depth grows with n); scaled from n=1000 \\
Python & DualPivotQuickSort & asc & 1.000.000 & ESTIMATED & same failure expected (recursion depth grows with n); scaled from n=1000 \\
Python & DualPivotQuickSort & desc & 10.000 & ERROR\_RecursionError & RecursionError at n=10000 (limit 1000); scaled from n=1000 \\
Python & DualPivotQuickSort & desc & 100.000 & ESTIMATED & same failure expected (recursion depth grows with n); scaled from n=1000 \\
Python & DualPivotQuickSort & desc & 1.000.000 & ESTIMATED & same failure expected (recursion depth grows with n); scaled from n=1000 \\
Python & InsertionSort & desc & 1.000.000 & ESTIMATED & predicted 22131s $>$ limit 300s; scaled from n=100000 \\
Python & InsertionSort & random & 1.000.000 & ESTIMATED & predicted 13718s $>$ limit 300s; scaled from n=100000 \\
\bottomrule
\end{tabular}}\par

\end{table}

\end{document}
